%% filename: amsbook-template.tex
%% version: 1.1
%% date: 2014/07/24
%%
%% American Mathematical Society
%% Technical Support
%% Publications Technical Group
%% 201 Charles Street
%% Providence, RI 02904
%% USA
%% tel: (401) 455-4080
%%      (800) 321-4267 (USA and Canada only)
%% fax: (401) 331-3842
%% email: tech-support@ams.org
%% 
%% Copyright 2006, 2008-2010, 2014 American Mathematical Society.
%% 
%% This work may be distributed and/or modified under the
%% conditions of the LaTeX Project Public License, either version 1.3c
%% of this license or (at your option) any later version.
%% The latest version of this license is in
%%   http://www.latex-project.org/lppl.txt
%% and version 1.3c or later is part of all distributions of LaTeX
%% version 2005/12/01 or later.
%% 
%% This work has the LPPL maintenance status `maintained'.
%% 
%% The Current Maintainer of this work is the American Mathematical
%% Society.
%%
%% ====================================================================

%    AMS-LaTeX v.2 driver file template for use with amsbook
%
%    Remove any commented or uncommented macros you do not use.

\documentclass{book}

%    For use when working on individual chapters
%\includeonly{}

%    For use when working on individual chapters
%\includeonly{}

%    Include referenced packages here.
\usepackage[left =.25in, right = .25in, top = 1in, bottom = 1in]{geometry} 
\usepackage{amsmath}
\usepackage{amsthm}
\usepackage{amscd}
\usepackage{amssymb}
\usepackage{setspace}
\usepackage{mathtools}
\usepackage{tikz}  
\usepackage{tikz-cd}
\usepackage{tkz-fct}
\usepackage{amsrefs}

\usepackage{pgfplots}
\pgfplotsset{compat=1.18}

\usepackage{environ}
\usepackage{tikz-cd} 
\usepackage{enumitem}
\usepackage{xcolor}   %May be necessary if you want to color links
%\usepackage{xr}

\usepackage{hyperref}
\hypersetup{
	colorlinks=true, %set true if you want colored links
	linktoc=all,     %set to all if you want both sections and subsections linked
	linkcolor=black,  %choose some color if you want links to stand out
	urlcolor=cyan
}
\usepackage[symbols,nogroupskip,sort=none]{glossaries-extra}

\pgfplotsset{every axis/.append style={
		axis x line=middle,    % put the x axis in the middle
		axis y line=middle,    % put the y axis in the middle
		axis line style={<->,color=black}, % arrows on the axis
		xlabel={$x$},          % default put x on x-axis
		ylabel={$y$},          % default put y on y-axis
}}


\theoremstyle{definition}
\newtheorem{definition}{Definition}[subsection]
\newtheorem{defn}[definition]{Definition}
\newtheorem{note}[definition]{Note}
\newtheorem{ax}[definition]{Axiom}
\newtheorem{thm}[definition]{Theorem}
\newtheorem{lem}[definition]{Lemma}
\newtheorem{prop}[definition]{Proposition}
\newtheorem{cor}[definition]{Corollary}
\newtheorem{conject}[definition]{Conjecture}
\newtheorem{ex}[definition]{Exercise}
\newtheorem{exmp}[definition]{Example}
\newtheorem{soln}[definition]{Solution}
\newtheorem{questn}[definition]{Question}

\setcounter{tocdepth}{3}

% hide proofs
\newif\ifhideproofs
%\hideproofstrue %uncomment to hide proofs
\ifhideproofs
\NewEnviron{hide}{}
\let\proof\hide
\let\endproof\endhide
\fi

% lower-case greek
\newcommand{\al}{\alpha}
\newcommand{\be}{\beta}
\newcommand{\gam}{\gamma}
\newcommand{\del}{\delta}
\newcommand{\ep}{\epsilon}
\newcommand{\ze}{\zeta} 
\renewcommand{\th}{\theta}
\newcommand{\kap}{\kappa} 
\newcommand{\lam}{\lambda}  
\newcommand{\sig}{\sigma} 
\newcommand{\omi}{\omicron}
\newcommand{\up}{\upsilon}
\newcommand{\om}{\omega}

% upper-case greek
\newcommand{\Gam}{\Gamma}
\newcommand{\Del}{\Delta}
\newcommand{\Lam}{\Lambda} 
\newcommand{\Sig}{\Sigma} 
\newcommand{\Om}{\Omega}


% text
\newcommand{\tA}{\text{A}}
\newcommand{\tB}{\text{B}}
\newcommand{\tC}{\text{C}}
\newcommand{\tD}{\text{D}}
\newcommand{\tE}{\text{E}}


% blackboard bold
\newcommand{\C}{\mathbb{C}}
\newcommand{\E}{\mathbb{E}}
\newcommand{\F}{\mathbb{F}}
\renewcommand{\H}{\mathbb{H}}
\newcommand{\K}{\mathbb{K}}
\newcommand{\N}{\mathbb{N}}
\renewcommand{\O}{\mathbb{O}}
\newcommand{\Q}{\mathbb{Q}}
\newcommand{\R}{\mathbb{R}}
\renewcommand{\S}{\mathbb{S}}
\newcommand{\T}{\mathbb{T}}
\newcommand{\V}{\mathbb{V}}
\newcommand{\Z}{\mathbb{Z}}

% math caligraphic
\newcommand{\MA}{\mathcal{A}}
\newcommand{\MB}{\mathcal{B}}
\newcommand{\MC}{\mathcal{C}}
\newcommand{\MD}{\mathcal{D}}
\newcommand{\ME}{\mathcal{E}}
\newcommand{\MF}{\mathcal{F}}
\newcommand{\MG}{\mathcal{G}}
\newcommand{\MH}{\mathcal{H}}
\newcommand{\MI}{\mathcal{I}}
\newcommand{\MJ}{\mathcal{J}}
\newcommand{\MK}{\mathcal{K}}
\newcommand{\ML}{\mathcal{L}}
\newcommand{\MM}{\mathcal{M}}
\newcommand{\MN}{\mathcal{N}}
\newcommand{\MO}{\mathcal{O}}
\newcommand{\MP}{\mathcal{P}}
\newcommand{\MQ}{\mathcal{Q}}
\renewcommand{\MR}{\mathcal{R}}
\newcommand{\MS}{\mathcal{S}}
\newcommand{\MT}{\mathcal{T}}
\newcommand{\MU}{\mathcal{U}}
\newcommand{\MV}{\mathcal{V}}
\newcommand{\MW}{\mathcal{W}}
\newcommand{\MX}{\mathcal{X}}
\newcommand{\MY}{\mathcal{Y}}
\newcommand{\MZ}{\mathcal{Z}}

% mathfrak
\newcommand{\MFX}{\mathfrak{X}}
\newcommand{\MFg}{\mathfrak{g}}
\newcommand{\MFh}{\mathfrak{h}}

% tilde 
\newcommand{\tMA}{\tilde{\MA}}
\newcommand{\tMB}{\tilde{\MB}}
\newcommand{\tU}{\tilde{U}}
\newcommand{\tV}{\tilde{V}}
\newcommand{\tphi}{\tilde{\phi}}
\newcommand{\tpsi}{\tilde{\psi}}
\newcommand{\tF}{\tilde{F}}

\newcommand{\iid}{\stackrel{iid}{\sim}}

\newcommand{\Cn}{\text{Cn}} % cone base shift functor

\newcommand{\cur}{\text{cur}} % curry operator





% label/reference
% internal label/reference
\newcommand{\lex}[1]{\label{ex:#1}}
\newcommand{\rex}[1]{Exercise \ref{ex:#1}}

\newcommand{\ld}[1]{\label{defn:#1}}
\newcommand{\rd}[1]{Definition \ref{defn:#1}}

\newcommand{\lnt}[1]{\label{note:#1}}
\newcommand{\rnt}[1]{Note \ref{note:#1}}

\newcommand{\lax}[1]{\label{ax:#1}}
\newcommand{\rax}[1]{Axiom \ref{ax:#1}}

\newcommand{\lfig}[1]{\label{fig:#1}}
\newcommand{\rfig}[1]{Figure \ref{fig:#1}}

\newcommand{\lsect}[1]{\label{sect:#1}}
\newcommand{\rsect}[1]{Section \ref{sect:#1}}
\newcommand{\rsectp}[1]{(Section \ref{sect:#1})}

% external reference
\newcommand{\extrex}[2]{Exercise \ref{#1-ex:#2}}

\newcommand{\extrd}[2]{Definition \ref{#1-defn:#2}}

\newcommand{\extrax}[2]{Axiom \ref{#1-ax:#2}}

\newcommand{\extrfig}[2]{Figure \ref{#1-fig:#2}}



% external documents (EDIT HERE)
%\externaldocument[analysis-]{"/home/carson/Desktop/Github/Mathematics/Introduction to Analysis/Introduction to Analysis.tex"}











% math operators
\DeclareMathOperator{\supp}{supp} % support
\DeclareMathOperator{\sgn}{sgn} % signum
\DeclareMathOperator{\spn}{span} % span

\DeclareMathOperator{\Obj}{Obj} % objects
\DeclareMathOperator{\dom}{dom} % domain
\DeclareMathOperator{\cod}{cod} % codomain

\DeclareMathOperator*{\ot}{\otimes} % otimes













\newcommand{\Oper}{\text{Op}}
\newcommand{\Ext}{\text{Ext}} % extensions
\newcommand{\Rstr}{\text{Rstr}} % restrictions
\newcommand{\OperExt}{\text{OpExt}}
\newcommand{\unb}{\text{unb}}























\DeclareMathOperator{\Eq}{Eq}

\DeclareMathOperator{\Sym}{Sym}
\DeclareMathOperator{\Alt}{Alt}









% linear algebra
\newcommand{\Dual}{\text{Dual}}
% ######################################################
% ######################################################
% ######################################################


































% ###################_________Sets_________###################
\newcommand{\Cover}{\text{Cover}} % cover

\DeclareMathOperator*{\prodrel}{\Pi^{\Rel}} % product of relations 
\DeclareMathOperator*{\timesrel}{\times^{\Rel}} % binary product of relations  

% Sets - Classes
\newcommand{\Set}{\text{Set}} % class of sets
\newcommand{\Rel}{\text{Rel}} % class of relations

% Sets - Categories
\newcommand{\bSet}{\text{\tbf{Set}}} % category of sets
\newcommand{\bRel}{\text{\tbf{Rel}}} % category of relations
% ######################################################
% ######################################################
% ######################################################
















% ###################_________Logic_________###################

\DeclareMathOperator{\Symb}{\text{Symb}} % set of symbols
\DeclareMathOperator{\RelSymb}{\text{RelSymb}} % set of relation symbols
\DeclareMathOperator{\FunSymb}{\text{FunSymb}} % set of function sympbols
\DeclareMathOperator{\ConsSymb}{\text{ConsSymb}}  % set of constant symbols
\DeclareMathOperator{\ConnSymb}{\text{ConnSymb}}  % set of connective symbols
\DeclareMathOperator{\Quan}{\text{Quan}}  % 
\DeclareMathOperator{\AuxSymb}{\text{AuxSymb}}  % set of auxiliary symbols
\DeclareMathOperator{\Term}{\text{Term}}  % set of terms
\DeclareMathOperator{\ElmTerm}{\text{ElmTerm}}  % elementary terms
\DeclareMathOperator{\AtomForm}{\text{AtomForm}} % atomic formulae

\DeclareMathOperator{\existsop}{\exists} % exists (operator)
% ######################################################
% ######################################################
% ######################################################























% ###################_________Topology_________###################

\DeclareMathOperator{\Homeo}{\text{Homeo}}

\DeclareMathOperator{\cl}{cl} % closure
\DeclareMathOperator{\Int}{Int} % interior
\DeclareMathOperator{\clset}{clSets} % closed supersets
\DeclareMathOperator{\Intset}{IntSets} % open subsets
\DeclareMathOperator*{\Punct}{\text{Punct}} 
\newcommand{\opn}[1]{{#1}^{\text{opn}}} % open adjective
\DeclareMathOperator{\FIP}{\text{FIP}} % finite intersection property

% Metric Spaces 
\newcommand{\Bnd}{\text{Bnd}} % set of bounded functions
\newcommand{\SemiNorm}{\text{SemiNorm}} % set of seminorms
\newcommand{\taulc}{\tau^{\text{lc}}} 


\newcommand{\Lip}{\text{Lip}}  % set of Lipschitz functions
\newcommand{\LipConst}{\text{LipConst}}  % set of Lipschitz constants of a function

\newcommand{\lip}{\text{lip}}  % Lipschitz constant of a function

\newcommand{\diln}{\text{diln}}  % dilation of a function 

\newcommand{\Alx}{\text{Alx}}

% Classes
\DeclareMathOperator*{\Top}{\text{Top}}
\DeclareMathOperator*{\TopEq}{\tbf{TopEq}}
\DeclareMathOperator*{\Frm}{\tbf{Frm}}
\DeclareMathOperator*{\Loc}{\tbf{Loc}}  

% Categories
\DeclareMathOperator*{\bTop}{\text{\tbf{Top}}} % category of topological spaces 
\DeclareMathOperator*{\bTopEq}{\text{\tbf{TopEq}}} % category of topological spaces with equivalence relations
\DeclareMathOperator*{\bFrm}{\text{\tbf{Frm}}} % category of frames
\DeclareMathOperator*{\bLoc}{\text{\tbf{Loc}}} % category of Locales
% ######################################################
% ######################################################
% ######################################################











% ###################_________Measure and Integration_________###################
\newcommand{\sigAlg}{\sig\text{-Alg}}

\newcommand{\Darb}{\text{Darb}} % Darboux integrable functions
\newcommand{\Riem}{\text{Riem}} % Riemann integrable functions
\newcommand{\IntSet}{\text{IntSet}}
\newcommand{\Tag}{\text{Tag}}
\newcommand{\mesh}{\text{mesh}}

\newcommand{\dt}{\, d t}
\newcommand{\ds}{\, d s}
\newcommand{\dx}{\, d x}

\newcommand{\ddel}{\, d \del}
\newcommand{\dmu}{\, d \mu}
\newcommand{\dnu}{\, d \nu}
\newcommand{\dlam}{\, d \lambda}
\newcommand{\dP}{\, d P}
\newcommand{\dQ}{\, d Q}
\newcommand{\dm}{\, d m}
\newcommand{\dsh}{\, d \#}


% Classes
\DeclareMathOperator{\Meas}{\text{Meas}}
\DeclareMathOperator{\TopMeas}{\text{TopMeas}}

\DeclareMathOperator{\Msr}{\text{Msr}}
\DeclareMathOperator{\Msrone}{\text{Msr}_1} 
\DeclareMathOperator{\MsrC}{\text{Msr}_{\C}} 

\DeclareMathOperator{\RadMsr}{\text{RadMsr}}
\DeclareMathOperator{\RadMsrone}{\text{RadMsr}_1}
\DeclareMathOperator{\RadMsrC}{\text{RadMsr}_{\C}}

\DeclareMathOperator{\TopMsr}{\text{TopMsr}}
\DeclareMathOperator{\TopMsrone}{\text{TopMsr}_1} 
\DeclareMathOperator{\TopMsrC}{\text{TopMsr}_{\C}} 

\DeclareMathOperator{\TopRadMsr}{\text{TopRadMsr}}
\DeclareMathOperator{\TopRadMsrone}{\text{TopRadMsr}_1}
\DeclareMathOperator{\TopRadMsrC}{\text{TopRadMsr}_{\C}} 


% Categories
\DeclareMathOperator{\bMeas}{\text{\tbf{Meas}}}
\DeclareMathOperator{\bTopMeas}{\text{\tbf{TopMeas}}}

\DeclareMathOperator{\bMsr}{\text{\tbf{Msr}}} 
\DeclareMathOperator{\bMsrone}{\text{\tbf{Msr}}_1} 
\DeclareMathOperator{\bMsrC}{\text{\tbf{Msr}}_{\C}} 

\DeclareMathOperator{\bRadMsr}{\text{\tbf{RadMsr}}}
\DeclareMathOperator{\bRadMsrone}{\text{\tbf{RadMsr}}_1}
\DeclareMathOperator{\bRadMsrC}{\text{\tbf{RadMsr}}_{\C}}

\DeclareMathOperator{\bTopMsr}{\text{\tbf{TopMsr}}} 
\DeclareMathOperator{\bTopMsrone}{\text{\tbf{TopMsr}}_1} 
\DeclareMathOperator{\bTopMsrC}{\text{\tbf{TopMsr}}_{\C}} 

\DeclareMathOperator{\bTopRadMsr}{\text{\tbf{TopRadMsr}}}
\DeclareMathOperator{\bTopRadMsrone}{\text{\tbf{TopRadMsr}}_1}
\DeclareMathOperator{\bTopRadMsrC}{\text{\tbf{TopRadMsr}}_{\C}} 
% ######################################################
% ######################################################
% ######################################################


































% ###################_________Probability_________###################

\newcommand{\Filtn}{\text{Filtrn}}
% ######################################################
% ######################################################
% ######################################################













































% ###################_________Geometry_________###################

\DeclareMathOperator{\Atlasinf}{\text{Atlasinf}} % collection of smooth atlases on a topological manifold
\DeclareMathOperator{\Diffeo}{\text{Diffeo}}


\newcommand{\germ}{\text{germ}} % germ of functions at at point
\newcommand{\germinf}{\text{germ}^{\infty}} % germ of functions at at point

\newcommand{\Paths}{\text{Paths}} % Paths of a manifold

\newcommand{\LocTriv}{\text{LocTriv}} % local trivializations
\newcommand{\LocTrivz}{\text{LocTriv}^0} % topological local trivializations
\newcommand{\LocTrivinf}{\text{LocTriv}^{\infty}} % smooth local trivializations

\newcommand{\Lie}{\text{Lie}} % Lie algebra of a Lie group


% Geom - Classes
\newcommand{\Manz}{\text{Man}^0} % class of topological manifolds
\newcommand{\Maninf}{\text{Man}^{\infty}}  % class of smooth manifolds with empty boundary
\newcommand{\ManBndinf}{\text{ManBnd}^{\infty}} % class of manifolds with nonempty boundary 

\newcommand{\LieGrp}{\text{LieGrp}} % class of Lie groups


% Geom - Categories
\DeclareMathOperator*{\bManz}{\text{\tbf{Man}}^0} % cat of topological manifolds
\DeclareMathOperator*{\bManinf}{\text{\tbf{Man}}^{\infty}} % cat of smooth manifolds with empty boundary
\DeclareMathOperator*{\bManBndinf}{\text{\tbf{ManBnd}}^{\infty}} % cat of manifolds with nonempty boundary 

\DeclareMathOperator*{\bBunz}{\text{\tbf{Bun}}^0} % cat of topological fiber bundles
\DeclareMathOperator*{\bBuninf}{\text{\tbf{Bun}}^{\infty}} % cat of smooth fiber bundles
\DeclareMathOperator*{\bVecBuninf}{\text{\tbf{VecBun}}^{\infty}} % cat of smooth vector bundles
\DeclareMathOperator*{\bPrinBuninf}{\text{\tbf{PrinBun}}^{\infty}} % cat of smooth principle bundles 

\DeclareMathOperator*{\bLieGrp}{\text{\tbf{LieGrp}}} 



\DeclareMathOperator{\Conn}{\text{Conn}}  % set of connections
% ######################################################
% ######################################################
% ######################################################

















% ###################_________Convexity_________###################

\DeclareMathOperator{\bal}{bal} % balanced subset generated by
\DeclareMathOperator{\cyc}{cyc} % 
\DeclareMathOperator{\cnv}{conv} % convex subset generate by


\DeclareMathOperator{\cnvspn}{\text{convspan}} 
\DeclareMathOperator{\extrm}{extrm} 


\DeclareMathOperator{\epi}{epi} % epigraph 
\DeclareMathOperator{\env}{env} % convex envelope

\newcommand{\Conv}{\text{Conv}} % convex subsets
\newcommand{\Bal}{\text{Bal}} % balanced subsets
\newcommand{\Aff}{\text{Aff}} % affine functions

\newcommand{\Cone}{\text{Cone}} % set of cones
\newcommand{\ConvCone}{\text{ConvCone}} % set of convex cones
\newcommand{\PtConvCone}{\text{PtConvCone}} % set of pointed convex cone
% ######################################################
% ######################################################
% ######################################################
 



\DeclareMathOperator{\codim}{codim} % codimension

\DeclareMathOperator{\Derivinf}{Deriv^{\infty}}

\DeclareMathOperator*{\argmax}{arg\,max}  
\DeclareMathOperator*{\argmin}{arg\,min}
\DeclareMathOperator{\diam}{\text{diam}} % diameter
\DeclareMathOperator{\rnk}{\text{rank}} % rank
\DeclareMathOperator{\tr}{\text{tr}} % trace
\DeclareMathOperator{\prj}{\text{proj}} % projection
\DeclareMathOperator{\nab}{\nabla} % nabla
\DeclareMathOperator{\diag}{\text{diag}} % diagonal
\DeclareMathOperator*{\ind}{\text{ind}} % 

\newcommand{\Part}{\text{Part}} 

\DeclareMathOperator{\Var}{\text{Var}}













































\DeclareMathOperator{\Subcat}{Subcat} % SubCat
\DeclareMathOperator{\SubcatGen}{SubcatGen} % subcategory generated by morphisms
\newcommand{\leqcat}{\leq_\text{Cat}} % Subcat partial order

\DeclareMathOperator{\Subquiv}{Subquiv} % Subquiv
\DeclareMathOperator{\SubquivGen}{\text{SubquivGen}} % subquiver generated by vertices/edges
\newcommand{\leqquiv}{\leq_\text{Quiv}} % Subquiver partial order


\DeclareMathOperator{\EmbSubmaninf}{\text{EmbSubman}^{\infty}} % Subquiv
\DeclareMathOperator{\EmbSubmaninfGen}{\text{EmbSubman}^{\infty}\text{Gen}} % subquiver generated by vertices/edges
\newcommand{\leqembmaninf}{\leq_{\text{Man}^{\infty}}} % Subquiver partial order

\DeclareMathOperator{\lmt}{lim}





























% ###################_________Algebra_________###################


% Groups, Rings, Fields, Modules, Algebras
\newcommand{\inv}{\text{inv}}
\newcommand{\mult}{\text{mult}}
\newcommand{\add}{\text{add}}
\newcommand{\scalmult}{\text{scalar\textunderscore mult}}
%\DeclareMathOperator{\smult}{\text{smult}}

\DeclareMathOperator{\End}{\text{End}} 
\DeclareMathOperator{\rep}{\text{Rep}} 
\DeclareMathOperator*{\Equi}{\text{Equiv}}
\DeclareMathOperator*{\Cong}{\text{Cong}}



% Classes
\newcommand{\Mon}{\text{Mon}} % class of monoids
\newcommand{\Grp}{\text{Grp}} % class of groups
\newcommand{\Semgrp}{\text{Semgrp}} % class of semigroups

\newcommand{\Ring}{\text{Ring}} % class of rings
\newcommand{\Field}{\text{Field}} % class of fields

\newcommand{\Mod}{\text{Mod}} % class of modules

\newcommand{\Vect}{\text{Vect}} % class of vector spaces
\newcommand{\VectR}{\text{Vect}_{\R}} % class of real vector spaces
\newcommand{\VectC}{\text{Vect}_{\C}} % class of complex vector spaces
\newcommand{\VectK}{\text{Vect}_{\K}} % class of K-vector spaces

\newcommand{\Alg}{\text{Alg}} % class of algebras
\newcommand{\AlgR}{\text{Alg}_{\R}} % class of real algebras
\newcommand{\AlgC}{\text{Alg}_{\C}} % class of complex algebras
\newcommand{\AlgK}{\text{Alg}_{\K}} % class of K-algebras



% Categories
\DeclareMathOperator*{\bMon}{\text{\tbf{Mon}}} % cat of monoids
\DeclareMathOperator*{\bGrp}{\text{\tbf{Grp}}} % cat of groups
\DeclareMathOperator*{\bSemgrp}{\text{\tbf{Semgrp}}} % cat of semigroups

\DeclareMathOperator*{\bRing}{\text{\tbf{Ring}}} % cat of rings
\DeclareMathOperator*{\bField}{\text{\tbf{Field}}} % cat of fields

\DeclareMathOperator*{\bMod}{\text{\tbf{Mod}}} % cat of modules

\DeclareMathOperator*{\bVect}{\text{\tbf{Vect}}} % cat of vector spaces
\DeclareMathOperator*{\bVectR}{\text{\tbf{Vect}}_{\R}} % cat of real vector spaces
\DeclareMathOperator*{\bVectC}{\text{\tbf{Vect}}_{\C}} % cat of complex vector spaces
\DeclareMathOperator*{\bVectK}{\text{\tbf{Vect}}_{\K}} % cat of K-vector spaces

\DeclareMathOperator*{\bAlg}{\text{\tbf{Alg}}} % cat of algebras
\DeclareMathOperator*{\bAlgR}{\text{\tbf{Alg}}_{\R}} % cat of real algebras
\DeclareMathOperator*{\bAlgC}{\text{\tbf{Alg}}_{\C}} % cat of complex algebras
\DeclareMathOperator*{\bAlgK}{\text{\tbf{Alg}}_{\K}} % cat of K-algebras

\DeclareMathOperator{\leqalg}{\text{Alg}} % cat of subalgebras 
% ######################################################
% ######################################################
% ######################################################





% Universal Algebra
\DeclareMathOperator*{\ar}{\text{arity}} % arity
\DeclareMathOperator*{\Sg}{\text{Sg}} 
\DeclareMathOperator{\uni}{Uni}
\DeclareMathOperator{\oper}{Oper}
\DeclareMathOperator{\indoper}{IndOper}
\DeclareMathOperator{\typ}{Type}
\DeclareMathOperator{\typs}{Types}
\DeclareMathOperator{\typfun}{TypeFun}
\DeclareMathOperator{\subu}{SubUni}
\DeclareMathOperator{\suba}{SubAlg}



























% ###################_________Order Theory_________###################


\DeclareMathOperator{\upar}{\uparrow}
\DeclareMathOperator{\doar}{\downarrow}
\DeclareMathOperator{\PosRefl}{PosRefl}

\DeclareMathOperator{\Mtn}{\text{Mtn}} % set of monotone maps
\DeclareMathOperator{\Antn}{\text{Antn}} % set of antitone maps

\DeclareMathOperator{\RAdj}{RAdj}
\DeclareMathOperator{\rad}{\text{rad}}
\DeclareMathOperator{\lad}{\text{lad}}
\DeclareMathOperator{\RQAdj}{RQAdj}
\DeclareMathOperator{\LAdj}{LAdj}
\DeclareMathOperator{\LQAdj}{LQAdj}

\DeclareMathOperator{\RConn}{RConn}
\DeclareMathOperator{\LConn}{LConn}

\DeclareMathOperator{\Ideals}{Ideals}
\DeclareMathOperator{\Filters}{Filters}
\DeclareMathOperator{\IdealGen}{IdealGen}
\DeclareMathOperator{\FilterGen}{FilterGen}


\DeclareMathOperator*{\prmj}{\text{prime}_{\vee}} % join prime
\DeclareMathOperator*{\prmm}{\text{prime}_{\wedge}} % meet prime
\DeclareMathOperator*{\irrj}{\text{irred}_{\vee}} % join irreducible
\DeclareMathOperator*{\irrm}{\text{irred}_{\wedge}} % meet irreducible

\DeclareMathOperator{\supSet}{supSet} % supremum set
\DeclareMathOperator{\infSet}{infSet} % infimum set

\DeclareMathOperator{\lbd}{lb}
\DeclareMathOperator{\ubd}{ub}


% Classes

% Categories
\DeclareMathOperator*{\bPos}{\text{\tbf{Pos}}} % cat of posets
























% ###################_________Category Theory________###################

\newcommand{\opp}{\text{op}} % opposite functor

\DeclareMathOperator{\id}{id} % identity
\DeclareMathOperator{\Hom}{Hom} % hom set
\DeclareMathOperator{\Endo}{End} % endomorphisms
\DeclareMathOperator{\Iso}{Iso} % ismorphisms
\DeclareMathOperator{\Aut}{Aut} % automorphisms
\DeclareMathOperator{\Mono}{Mono} % monomorphisms
\DeclareMathOperator{\ExtMono}{ExtMono} % extremal monomorphisms
\DeclareMathOperator{\RegMono}{RegMono} % regular monomorphisms
\DeclareMathOperator{\Epi}{Epi} % epimorphisms
\DeclareMathOperator{\ExtEpi}{ExtEpi} % extremal monomorphisms
\DeclareMathOperator{\RegEpi}{RegEpi} % regular monomorphisms

%\DeclareMathOperator{\Term}{Term} % terminal objects in a category
\DeclareMathOperator{\Coterm}{Coterm} % coterminal objects in a category

\DeclareMathOperator{\Prod}{Prod} % product objects in a category
\DeclareMathOperator{\Coprod}{Coprod} % coproduct object in a category

\DeclareMathOperator{\Exp}{Exp} % exponential objects in a category
\DeclareMathOperator{\Coexp}{Coexp} % coexponential objects in a category

\DeclareMathOperator{\Tmnl}{Term} % terminal objects in a category
\DeclareMathOperator{\Init}{Init} % initial objects in a category

\DeclareMathOperator{\Eqzr}{Equi} % equilizer objects in a category
\DeclareMathOperator{\Coeqzr}{Coequi} % coequilizer objects in a category



\DeclareMathOperator{\bMono}{\tbf{Mono}} % subcategory restricted to monomorphisms
\DeclareMathOperator{\bEpi}{\tbf{Epi}} % subcategory restricted to monomorphisms

\DeclareMathOperator{\Sub}{Sub} % subobject proset

\DeclareMathOperator{\colim}{colim} % colimit of diagram
\DeclareMathOperator{\limSet}{limSet} % limit set of diagram
\DeclareMathOperator{\colimSet}{colimSet} % colimit set of diagram
\DeclareMathOperator{\limChoice}{limChoice} % limit choice function 
\DeclareMathOperator{\colimChoice}{colimChoice} % colimit choice function 

% Cateogries
\DeclareMathOperator{\bArr}{\textbf{Arr}} % arrow category
\DeclareMathOperator*{\bCat}{\text{\tbf{Cat}}} % category of small categories
\DeclareMathOperator*{\bCone}{\text{\tbf{Cone}}} % category of small cones
\DeclareMathOperator*{\bCocone}{\text{\tbf{Cocone}}} % category of small cocones


































\newcommand{\bQuiv}{\text{\tbf{Quiv}}}


\newcommand{\Factors}{\text{Factors}}
\newcommand{\Factor}{\text{Factor}}

\newcommand{\LFactors}{\text{LFactors}}
\newcommand{\LFactor}{\text{LFactor}}


\newcommand{\RFactors}{\text{RFactors}}
\newcommand{\RFactor}{\text{RFactor}}


\DeclareMathOperator*{\bBanAlg}{\text{\tbf{BanAlg}}}



\DeclareMathOperator*{\bNorm}{\text{\tbf{Norm}}} 
\DeclareMathOperator*{\bBan}{\text{\tbf{Ban}}} 
\DeclareMathOperator*{\bHilb}{\text{\tbf{Hilb}}} 
\DeclareMathOperator*{\bProb}{\text{\tbf{Prob}}} 

\DeclareMathOperator*{\bRep}{\text{\tbf{Rep}}} 
\DeclareMathOperator*{\bURep}{\text{\tbf{URep}}}





\DeclareMathOperator*{\bTopField}{\text{\tbf{TopField}}} 
\DeclareMathOperator*{\bTopMon}{\text{\tbf{TopMon}}} 
\DeclareMathOperator*{\bTopGrp}{\text{\tbf{TopGrp}}}
\DeclareMathOperator*{\bTopVect}{\text{\tbf{TopVect}}} 










%\newcommand{\TopMeas}{\text{TopMeas}}
%\newcommand{\Msrpos}{\text{Msr}_{+}}
%\newcommand{\TopMsrpos}{\text{TopMsr}_{+}}
%\newcommand{\TopRadMsrpos}{\text{TopRadMsr}_{+}}
%\newcommand{\TopRadMsrone}{\text{TopRadMsr}_{1}}
%\newcommand{\MsrC}{\text{Msr}_{\C}} 
%\newcommand{\TopMsrC}{\text{TopMsr}_{\C}} 
%\newcommand{\TopRadMsrC}{\text{TopRadMsr}_{\C}} 
%\newcommand{\Maninf}{\text{Man}^{\infty}} 
%\newcommand{\ManBndinf}{\text{ManBnd}^{\infty}} 
%\newcommand{\Man0}{\text{Man}^{0}}
%\newcommand{\Buninf}{\text{Bun}^{\infty}} 
%\newcommand{\VecBuninf}{\text{VecBun}^{\infty}} 
%
%\newcommand{\Mon}{\text{Mon}} 
%\newcommand{\Grp}{\text{Grp}}
%\newcommand{\Semgrp}{\text{Semgrp}}
%\newcommand{\LieGrp}{\text{LieGrp}} 
%
%\newcommand{\Field}{\text{Field}} 
%
%\newcommand{\Vect}{\text{Vect}}
%\newcommand{\Mod}{\text{Mod}}
%\newcommand{\Alg}{\text{Alg}} 
%
%\newcommand{\Rep}{\text{Rep}} 
%\newcommand{\URep}{\text{URep}}
%
%\newcommand{\bNorm}{\text{Norm}} 
%\newcommand{\Ban}{\text{Ban}} 
%\newcommand{\Hilb}{\text{Hilb}} 
%\newcommand{\BanAlg}{\text{BanAlg}}
%
%\newcommand{\Prob}{\text{Prob}} 
%\newcommand{\PrinBuninf}{\text{PrinBun}^{\infty}}
%
%\newcommand{\TopField}{\text{TopField}} 
%\newcommand{\TopMon}{\text{TopMon}}
%\newcommand{\TopGrp}{\text{TopGrp}}
%\newcommand{\TopVect}{\text{TopVect}} 
%\newcommand{\TopEq}{\text{TopEq}}
%\newcommand{\Frm}{\text{Frm}}
%\newcommand{\Loc}{\text{Loc}}  
%
%\newcommand{\VectR}{\text{Vect}_{\R}}
%\newcommand{\VectC}{\text{Vect}_{\C}} 
%\newcommand{\VectK}{\text{Vect}_{\K}}
%
%\newcommand{\Cat}{\text{Cat}
%\newcommand{\Cone}{\text{Cone}
%\newcommand{\Cocone}{\text{Cocone}


\newcommand{\0}{\mbf{0}}
\newcommand{\1}{\mbf{1}}
\newcommand{\2}{\mbf{2}}







\DeclareMathOperator*{\Fac}{\text{Factor}}




% notation
\renewcommand{\r}{\rangle}
\renewcommand{\l}{\langle}
\renewcommand{\div}{\text{div}}
\renewcommand{\Re}{\text{Re} \,}
\renewcommand{\Im}{\text{Im} \,}
\newcommand{\Img}{\text{Img} \,}
\newcommand{\grad}{\text{grad}}





% cat theory notation
\newcommand{\pred}{\text{pred}}
\newcommand{\comp}{\text{comp}}


\newcommand{\Bool}{\text{Bool}}


\newcommand{\tbf}[1]{\textbf{#1}}
\newcommand{\tcb}[1]{\textcolor{blue}{#1}}
\newcommand{\tcr}[1]{\textcolor{red}{#1}}
\newcommand{\mbf}[1]{\mathbf{#1}}


\newcommand{\ol}[1]{\overline{#1}}
\newcommand{\ub}[1]{\underbar{#1}}
\newcommand{\tl}[1]{\tilde{#1}}
\newcommand{\p}{\partial}
\newcommand{\Tn}[1]{T^{r_{#1}}_{s_{#1}}(V)}
\newcommand{\Tnp}{T^{r_1 + r_2}_{s_1 + s_2}(V)}
\newcommand{\Perm}{\text{Perm}}
\newcommand{\wh}[1]{\widehat{#1}}
\newcommand{\wt}[1]{\widetilde{#1}}
\newcommand{\defeq}{\vcentcolon=}

\newcommand{\KosConn}{\text{KosConn}}

\newcommand{\tri}{\triangle}
\newcommand{\trl}{\triangleleft}
\newcommand{\trr}{\triangleright}
\newcommand{\loz}{\lozenge}
\newcommand{\dmd}{\diamond}
\newcommand{\alg}{\text{alg}}
\newcommand{\Triv}{\text{Triv}}
\newcommand{\Der}{\text{Der}}


\DeclareMathOperator{\conj}{conj}
\DeclareMathOperator{\conjdom}{\text{dom}_{\text{conj}}}
\DeclareMathOperator{\biconj}{\text{biconj}}
\DeclareMathOperator{\biconjdom}{\text{dom}_{\text{biconj}}}



\newcommand{\lcm}{\text{lcm}}
\newcommand{\Imax}{\MI_{\text{max}}}


\DeclareMathOperator{\Rl}{\text{Re}}
\DeclareMathOperator{\Imn}{\text{Imn}}

\newcommand{\veq}{\rotatebox{90}{$\,=$}}
\newcommand{\dtbeq}{\rotatebox{-45}{$\,=$}}
\newcommand{\dbteq}{\rotatebox{45}{$\,=$}}




\newcommand{\eval}{\text{eval}}
\newcommand{\coeval}{\text{coeval}}




% limits



\newcommand{\limfn}{\liminf \limits_{n \rightarrow \infty}}
\newcommand{\limpn}{\limsup \limits_{n \rightarrow \infty}}

\newcommand{\limn}{\lim \limits_{n \rightarrow \infty}}
\newcommand{\limx}[1]{\lim \limits_{x \rightarrow #1}}
\newcommand{\limt}[1]{\lim \limits_{t \rightarrow #1}}
\newcommand{\limh}[1]{\lim \limits_{h \rightarrow #1}}

\newcommand{\limneq}[2]{\lim\limits_{\substack{#1 \rightarrow #2 \\ #1 \neq #2}}}
\newcommand{\limneqx}[1]{\lim\limits_{\substack{x \rightarrow #1 \\ x \neq #1}}}
\newcommand{\limneqt}[1]{\lim\limits_{\substack{t \rightarrow #1 \\ x \neq #1}}}

\newcommand{\limfneq}[2]{\liminf\limits_{\substack{{#1} \rightarrow #2 \\ #1 \neq #2}}}
\newcommand{\limfneqx}[1]{\liminf\limits_{\substack{x \rightarrow #1 \\ x \neq #1}}}

\newcommand{\convt}[1]{\xrightarrow{\text{#1}}}
\newcommand{\conv}[1]{\xrightarrow{#1}} 
\newcommand{\seq}[2]{(#1_{#2})_{#2 \in \N}}

% intervals
\newcommand{\Rext}{[-\infty,\infty]}


\newcommand{\RG}{[0,\infty]}
\newcommand{\Rg}{[0,\infty)}
\newcommand{\Rgp}{(0,\infty)}
\newcommand{\Ru}{(\infty, \infty]}
\newcommand{\Rd}{[\infty, \infty)}
\newcommand{\ui}{[0,1]}

\newcommand{\RE}{\ol{\R}}
\newcommand{\REp}{\ol{\R}_+}

\newcommand{\Rp}{\R_+}







% abreviations 
\newcommand{\lsc}{lower semicontinuous}

% misc
\newcommand{\as}[1]{\overset{#1}{\sim}}
\newcommand{\astx}[1]{\overset{\text{#1}}{\sim}}
\newcommand{\io}{\text{ i.o.}}
%\newcommand{\ev}{\text{ ev.}}
\newcommand{\Ll}{L^1_{\text{loc}}(\R^n)}

\newcommand{\loc}{\text{loc}}
\newcommand{\BV}{\text{BV}}
\newcommand{\NBV}{\text{NBV}}
\newcommand{\TV}{\text{TV}}

\newcommand{\op}[1]{{#1}^{\text{op}}}
\newcommand{\lop}{\leq^*}


% Glossary - Notation
\glsxtrnewsymbol[description={finite measures on $(X, \MA)$}]{n000001}{$\MM_+(X, \MA)$}
\glsxtrnewsymbol[description={velocity}]{v}{\ensuremath{v}}


\makeindex

\begin{document}
	
	\frontmatter
	
	\title{Introduction to Mathematics}
	
	%    Remove any unused author tags.
	
	%    author one information
	\author{Carson James}
	\thanks{}
	
	\date{}
	
	\maketitle
	
	%    Dedication.  If the dedication is longer than a line or two,
	%    remove the centering instructions and the line break.
	%\cleardoublepage
	%\thispagestyle{empty}
	%\vspace*{13.5pc}
	%\begin{center}
	%  Dedication text (use \\[2pt] for line break if necessary)
	%\end{center}
	%\cleardoublepage
	
	%    Change page number to 6 if a dedication is present.
	\setcounter{page}{4}
	
	\tableofcontents
	\printunsrtglossary[type=symbols,style=long,title={Notation}]
	
	%    Include unnumbered chapters (preface, acknowledgments, etc.) here.
	%\include{}
	
	\mainmatter
	%    Include main chapters here.
	%\include{}
	
	\chapter*{Preface}
	\addcontentsline{toc}{chapter}{Preface}
	
	\begin{flushleft}
		\href{https://creativecommons.org/licenses/by-nc-sa/4.0/legalcode.txt}{cc-by-nc-sa}
	\end{flushleft}
	
	\newpage
	
	
	
	
	
	
	
	
	
	
	
	
	
	
	
	
	
	
	
	
	
	
	
	
	
	
	
	
	
	
	
	
	\newpage
	\part{Logic and Set Theory}
	\section{To Do}
	\begin{itemize}
		\item need to merge the set theory notes with this intro to math. It might as well cover all the stuff related to set theory, just with more prose. 
	\end{itemize}
	
	\chapter{Introduction to Logic}
	
	\subsection{Quantification}
	
	\begin{itemize}
		\item discuss propositional logic and how it can represent statements and how to model it with computer say. Talk about conjunction (and), disjunction (or), implication (if,then), exlusive disjunction (either or) \tcr{(should discuss the controversy about `either or'. Just note that human language is often imprecise and establish a convention that `or' is to be taken in the sense of a choice of soup or salad at a buffet while `either or' is to be taken in the sense of a choice of soup or salad at a restaurant)}.    
		\item Motivate quantification and type theory as a way to address the need for a conjucntion/disjunction for infintely many statements while propositional logic only accomodates conjunction of finitely many statements. For example we would like to be able to represent/reason with statements like ``there is a positive integer that is larger than $2$" or ``each integer is positive". We can intuit these types of statements in a few ways: 
		\begin{itemize}
			\item One way is with infinitely many dconjunctions/isjuctions: We can interpret the statement ``there is a positive integer that is larger than $2$" as the infinte disjunction ``$1 > 2$ or $2 > 2$ or $3 > 2$ or $\ldots$". Since this disjunction does not exist in propositional logic, we can just define an extension of propositional logic in which the infinte disjunction exists and is a well-defined operation. This is analogous to how the positive real numbers $\Rg$ does not have a natural infinte summation operation. We resolve both of these issues by extending the objects under consideration, statements $\phi$, (elements of $\Rg$ in the analogy), to a larger collection of objects, indexed statements $(\phi(x))_{x \in \MI}$, such as ``$(x > 2)_{x \in \N}$", where we indentify statements from the original collection such as ``$5$ is odd" with the indexed family $(\phi(x))_{x \in \N}$ where $\phi(x) \defeq \text{``$5$ is odd"}$ if $x = 5$ and $\phi(x) \defeq \perp$ if $x \neq 5$ ($(a_n)_{n \in \N} \in \RG^{\N}$, i.e. sequences of nonegative real numbers in the analogy where we identify $a \in \Rg$ with $(a_n)_{n \in \N}$ defined by $a_1 \defeq a$ and $a_n \defeq 0$). Then we define $\text{Value}[\existsop\limits_{x \in \MI} \phi(x)] \defeq \sup\limits_{x \in \MI} \text{Value}(\phi(x))$. 
			\item Another is as an implication: We can interpret the statement ``each integer is positive" as ``$a \in \Z \implies a >0$. In other words, if the type of $a$ (whatever $a$ that means) is $\Z$, then $a > 0$. To prove this statement, we would need to start with an object $a$ of type $\Z$, then provide a proof that $a > 0$. 
		\end{itemize}
		These both suffer from ambiguities, as we don't know what $\Z$ (or sets/sequences/etc in general) nor ``type" are. To get around this we create a formal system in which we assume ways to manipulate strings of characters according to specified rules which are aligned with propositional calculus (i.e. something we can run on a computer). try to capture the essence of quantification or type with these rules. 
		\item For quantification and set theory, .  
		\item Discuss quantification from a set theory perspective where everything is a set and we quantify over sets. For example, we would represent `` `is even', we would write `$\forall x, x$" 
		\item \tcr{discuss phrasing $\forall$ and $\exists$ in terms of a hypothetical sequence of transactions between two people, so that $f: \R \rightarrow \R$ is continuous at $x \in \R$ if $\forall \ep >0, \exists \del >0, \forall y \in \R, |x-y|< \del \implies |f(x) - f(y)| < \ep$ can be rephrased as ``If you give me an $\ep>0$ (doesn't matter which one you give me), I can then give you a $\del >0$ (which may depend of $\ep$) such that (i.e. that satisfies/has the property that) if you give me a $y \in \R$ (doesn't matter which one), we have that (I can show you that) if $|x-y|<\del$, then $|f(x) - f(y)| < \ep$."}
		\item \tcr{discuss the connection between restricted quantification and unrestricted quantification with implication. For example, ``$\forall x \in \R$, $x^2 > 0$" can be rephrased as ``$\forall x, x \in \R \implies x^2 > 0$". Point out that some sources claim that restricted quantification is equivalent to implication. For example ``$\forall x \in \R$, $x$" are equivalent and why we use quantification in the first place. For example } The reason we 
	\end{itemize}
	
	
	
	
	
	
	
	
	
	
	
	
	
	
	
	
	
	
	
	
	
	
	
	
	
	
	
	
	
	
	
	
	
	
	
	
	
	
	
	
	
	
	
	
	
	
	
	
	
	\newpage
	
	\chapter{Introduction to Set Theory}
	
	\section{Sets}	
	
	In this section, we provide an exposure to a basic framework, called set theory. This framework provides a unifying language which allows us to discuss many objects of interest in mathematics by providing the ability to 
	\begin{enumerate}
		\item more precisely define objects of interest, 
		\item more precisely formulate and prove statements about those objects
		\item perform (1) and (2) in a way which is aligned with our intuition.
	\end{enumerate}
	At an intuitive level, the main idea of a \tbf{set} is that of a collection of distinct objects, which we refer to as \tbf{elements} or \tbf{members}. When denoting a set, it is customary to list the elements of the set, separated by a comma and inside of a pair of curly brackets. Some examples of sets include the following:
	\begin{itemize}
		\item the set of letters of the English alphabet: $\{a, b, \ldots, z\}$, a set with 26 elements
		\item the set of inhabitants of a specific household: $\{\text{``Marielle"}, \text{``Fred"}, \text{``Dexter"}\}$, a set with 3 elements
		\item the set of nonnegative even integers: $\{0, 2, \ldots \}$, a set with infintely many elements
	\end{itemize}
	\begin{note}
		Since the elements of a set must be distinct, we would consider, for instance, the sets $\{a, b, c\}$ and $\{a, a, b, c\}$ to represent the same set. 
	\end{note}
	When expressing that a given object is a member of set, it is customary to use the \tbf{membership symbol}, denoted $\in$. Some examples include the following:
	\begin{itemize}
		\item $2 \in \{1, 2, 3, 4\}$
		\item $\text{``Earth"} \in \{\text{``Mercury"}, \text{``Venus"}, \text{``Earth"}, \text{``Mars"}\}$.
	\end{itemize}
	\begin{note}
		We note that this introduction to set theory is not formal, but intuitve. As such, we make no attempt to precisely define the notion of a set or the notion of membership to a set. 
	\end{note}
	At this point in our discussion, it might seem that any advantage conferred by the language of sets in studying mathematics remains to be seen. Indeed, more time and experience may be required to fully appreciate the power of this framework, but for now we summarize why set theory is relevant in mathematics with the following mantra: 
	$$\text{In mathematics, everything is a set.}$$
	There is a similar perspective adopted in the context of Linux operating systems where every object is a file or in the context of object oriented programming where every object is a class. 
	
	If one already has familiarity with various objects in mathematics such as various numbers, functions, and shapes, they might wonder how these objects could be considered as a set, i.e. a list of elements. For example, which elements $a_1, a_2, \ldots$, which are themselves sets (because everything is a set), would allow one to express the integer $2$ as the set $\{a_1, a_2, \ldots\}$? At first glance, we might attempt to represent $2$ with an arbitrary set containing two elements, say $\{\text{``cow"}, \text{``chicken"}\}$. The intuition here might be that the number of elements reflects our notion of the ``quantity" or ``magnitude" we associate to the number $2$. However, upon further reflection, such a set might not satisfy other notions we associate to the number $2$. 
	
	For example we typically think of the number $2$ as having a unique ``ranking" or ``place" in the number line. That is, the number $2$ is the integer that preceeds the number $3$. So we would need to not only define a set representing the number $3$, but also a set representing the notion of subsequence that would align with our intuition that $3$ is subsequent to $2$. 
	
	In a different vein, one might also wish for our representation of the number $2$ to agree with the idea that the sum of $2$ and $4$ is $6$. The issue here is how to describe the notion of integer addition. After some thought, one might realize that the idea of adding two integers is basically that of a function, i.e. an object which eats two integer inputs and spits out an integer output. 
	
	It might seem tricky to be able to represent all the nonnegative integers, the notion of integer subsequence and the notion of integer addition by sets in a way that aligns with our intution of those objects, relations and operations. With our limited toolset, this would indeed be quite challenging, but after a few more ideas, this will become a relatively simple goal. For now, we focus on providing the necessary calculus of sets to be able to do so. 
	
	The first two relations between sets that we define pertain to when one set is a subset of another and when one set is equal to another. To motivate this, we observe that each element of $\{a,b\}$ is also an element of $\{a, b, c\}$. Our intuition dictates that the former is in some sense ``less than or equal to" the latter. However, the same could not be said of the converse. Therefore, the two sets should not be considered equal.
	\begin{defn}
		Let $A$ and $B$ be sets. Then 
		\begin{itemize}
			\item $A$ is said to be a \tbf{subset of $B$} if for each set $x$, $x \in A$ implies that $x \in B$. \\
			($\text{$A$ is a subset of $B$} \iff \forall x (x \in A \implies x \in B)$)
			\item $A$ is said to be \tbf{equal to $B$} if $A \subset B$ and $B \subset A$. 
			\item $A$ is said to be \tbf{empty} if for each set $x$, $x \not \in A$. \\
			($\text{$A$ is empty} \iff \forall x ( \neg(x \in A))$)
		\end{itemize}
	\end{defn}
	
	\begin{note}\
		\begin{itemize}
			\item When defining a term, it is conventional to use an ``if", but really the relationship is an ``iff".
			\item For convenience, we will typically write $x \not \in A$ in place of $\neg(x \in A)$.
		\end{itemize}
	\end{note}
	
	The following exercise shows us that there is at most one set that is empty.
	\begin{ex}
		Let $A, B$ be sets. If $A$ is empty and $B$ is empty, then $A= B$. 
	\end{ex}
	
	\begin{proof}
		Suppose that $A$ is empty and $B$ is empty. Let $x$ be a set. 
		\begin{itemize}
			\item Since $A$ is empty, $x \not \in A$. Therefore $x \in A$ implies that $x \in B$. Since $x$ is an arbitrary set, we have that for each set $x$, $x \in A$ implies that $x \in B$. Hence $A \subset B$.
			\item Similarly $B \subset A$
		\end{itemize}
		Since $A \subset B$ and $B \subset A$, we have that $A = B$.
	\end{proof}
	
	Even though there can be at most one set that is empty, we are not guaranteed the existence of such a set. Since its existence is quite useful and intuitively quite reasonable, we will assume such a set does in fact exist. 
	\begin{ax}
		There exists a set $\varnothing$, referred to as the \tbf{empty set}, such that for each set $x$, $x \not \in \varnothing$.
	\end{ax}
	
	
	\begin{ex}
		Let $A$ be a set. Then $\varnothing \subset A$. 
	\end{ex}
	
	\begin{proof}
		Let $x$ be a set. Then $x \not \in \varnothing$. Hence $x \in \varnothing$ implies that $x \in A$. Since $x$ is an arbitrary set, we have that for each set $x$, $x \in \varnothing$ implies that $x \in A$. Thus $\varnothing \subset A$.
	\end{proof}
	
	\tcr{give intuition for union, intersection, complement, show we can define union in terms of intersection and complement, call back to demorgans laws for logic}
	\begin{defn}
		Let $A, B$ be sets. We define 
		intersection, union, empty set
	\end{defn}
	
	
	
	\tcr{discuss and introduce/make pictures of venn diagrams}
	
	
	\tcr{discuss how definitions by convention use an if, but it should be treated as an iff. The if makes sense because if we say that a dog is ``blarge" if it is black and weighs more than 15 kg, if we had put an iff in this definition, then the we would be saying that if a dog is blarge, then it is black and weighs more than 15 kgs, which wouldn really make sense in a definition because we have not defined previously and do not know what it is blarge, so it is a language thing. but really we are defining a set. $x \in \{\text{blarge dogs}\} \iff x \in \{\text{black dogs}\} \text{and} \{\text{dogs above 15 kg}\}$} i.e. $\{\text{blarge dogs}\} = \{\text{black dogs}\} \cap \{\text{dogs above 15 kg}\}$.
	
	
	
	
	discuss how we can define natural numbers
	
	discuss how we define $(a,b) = \{\{a\}, \{a,b\}\}$ or similar
	
	\begin{defn}
		\tcr{products $A \times B$}
	\end{defn}
	
	\begin{ex}
		content...
	\end{ex}
	
	\begin{ex}
		
	\end{ex}
	
	
	
	
	
	
	
	
	
	
	
	
	
	
	
	
	
	
	
	
	
	
	
	
	
	
	
	
	
	
	
	
	
	
	
	
	
	
	
	
	
	
	
	
	
	
	
	
	
	
	
	
	
	
	
	
	\newpage
	\section{Relations}
	
	\tcr{FINISH!!!}
	
	
	
	
	
	
	
	
	
	
	
	
	
	
	
	
	
	
	
	
	
	
	
	
	
	
	
	
	
	
	
	
	
	
	
	
	
	
	
	
	
	
	
	
	
	
	
	
	
	
	
	
	\newpage
	\section{Functions}
	
	\subsection{Introduction}
	
	\begin{defn}
		\tcr{define function}, 
		Let $X,Y$ be sets. 
		\begin{itemize}
			\item Let $f \subset \MR(A,B)$. Then $f$ is said to be a \tbf{function from $X$ to $Y$}, denoted $f:X \rightarrow Y$, if 
			\begin{enumerate}
				\item for each $x \in X$, there exists $y \in Y$ such that $(x,y) \in f$,
				\item for each $x_1, x_2 \in X$, if $x_1 = x_2$, then $f(x_1) = f(x_2)$. 
			\end{enumerate}
			\item We define $Y^X \defeq \{f: f:X \rightarrow Y\}$.
		\end{itemize}
		\tcr{later in section on products, show that $Y^X = \prod\limits_{x \in X} A$}
	\end{defn}
	
	\begin{ex}
		Let $X$ be a set. If $X \neq \varnothing$, then $\varnothing^{X} = \varnothing$. 
	\end{ex}
	
	\begin{proof}
		For the sake of contradiction, suppose that $X \neq \varnothing$ and $\varnothing^{X} \neq \varnothing$. Then there exists $f:X \rightarrow \varnothing$. Since $X \neq \varnothing$, there exists $x \in X$. Since $f$ is a function, there exists $y \in \varnothing$ such that $f(x) = y$. This is a contradiction. Hence $\varnothing^{X} = \varnothing$.
	\end{proof}
	
	\begin{defn}
		Let $X,Y$ be sets and $f:X \rightarrow Y$. Then $f$ is said to be 
		\begin{itemize}
			\item \tbf{injective} or an \tbf{injection} if for each $x_1, x_2 \in X$, $f(x_1) = f(x_2)$ implies that $x_1 = x_2$.
			\item \tbf{surjective} or a \tbf{surjection} if for each $y \in Y$, there exists $x \in X$ such that $f(x) = y$
			\item \tbf{bijective} a \tbf{bijection} if $f$ is injective and $f$ is surjective.
		\end{itemize}
	\end{defn}
	
	\begin{ex}
		Let $X,Y$ be sets and $f:X \rightarrow Y$. Then 
		\begin{enumerate}
			\item $f$ is injective iff for each $y \in f(X)$, there exists $x \in X$ such that $f^{-1}(\{y\}) = \{x\}$,
			\item $f$ is surjective iff $f(X) = Y$.
		\end{enumerate}
	\end{ex}
	
	\begin{proof}\
		\begin{enumerate}
			\item 
			\begin{itemize}
				\item $(\implies)$: \\
				Supose that $f$ is injective. Let $y \in f(X)$. Then there exists $x \in X$ such that $f(x) = y$. 
				\begin{itemize}
					\item Since 
					\begin{align*}
						f(x)
						& = y \\
						& \in \{y\},
					\end{align*}
					we have that $x \in f^{-1}(\{y\})$. Therefore $\{x\} \subset f^{-1}(\{y\})$. 
					\item Let $x' \in f^{-1}(\{y\})$. By definition \tcr{(ref def)}, $f(x') \in \{y\}$. \rex{} \tcr{(ref ex)} implies that $f(x') = y$. Therefore
					\begin{align*}
						f(x')
						& = y \\
						& = f(x).
					\end{align*}
					Since $f$ is injective, $x' = x$. Hence $x' \in \{x\}$. Since $x' \in f^{-1}(\{y\})$ is arbitrary, we have that for each $x' \in f^{-1}(\{y\})$, $x' \in \{x\}$. Therefore $f^{-1}(\{y\}) \subset \{x\}$. 
				\end{itemize}
				Since $\{x\} \subset f^{-1}(\{y\})$ and $f^{-1}(\{y\}) \subset \{x\}$, we have that $f^{-1}(\{y\}) = \{x\}$. Since $y \in f(X)$ is arbitrary, we have that for each $y \in f(X)$, there exists $x \in X$ such that $f^{-1}(\{y\}) = \{x\}$. 
				\item $(\impliedby)$: \\
				Suppose that for each $y \in f(X)$, there exists $x \in X$ such that $f^{-1}(\{y\}) = \{x\}$. Let $x_1, x_2 \in X$. Suppose that $f(x_1) = f(x_2)$. Define $y \in f(X)$ by $y \defeq f(x_1)$. By assumption, there exists $x \in X$ such that $f^{-1}(\{y\}) = \{x\}$. Since 
				\begin{align*}
					f(x_1)
					& = y \\
					& \in \{y\},
				\end{align*}
				we have that $x_1 \in f^{-1}(\{y\})$. Therefore 
				\begin{align*}
					x_1 
					& \in f^{-1}(\{y\}) \\
					& = \{x\}.
				\end{align*}
				Similarly, $x_2 \in \{x\}$. \rex{} \tcr{(ref ex)} then implies that 
				\begin{align*}
					x_1
					& = x \\
					& = x_1.
				\end{align*}
				Since $x_1, x_2 \in X$ with $f(x_1) = f(x_2)$ are arbitrary, we have that for each $x_1, x_2 \in X$, $f(x_1) = f(x_2)$ implies that $x_1 = x_2$. Thus $f$ is injective. 
			\end{itemize}
			\item 
			\begin{itemize}
				\item $(\implies)$: \\
				Supose that $f$ is surjective. 
				\begin{itemize}
					\item \rd{} \tcr{(ref def)} implies that $f(X) \subset Y$. 
					\item Let $y \in Y$. Since $f$ is surjective, there exists $x \in X$ such that $f(x) = y$. Then 
					\begin{align*}
						y
						& = f(x) \\
						& \in f(X).
					\end{align*}
					Since $y \in Y$ is arbitrary, we have that for each $y \in Y$, $y \in f(X)$. Hence $Y \subset f(X)$. 
				\end{itemize}
				Since $f(X) \subset Y$ and $Y \subset f(X)$, we have that $f(X) = Y$.
				\item $(\impliedby)$: \\
				Suppose that $f(X) = Y$. Let $y \in Y$. Since $f(X) = Y$, we have that $y \in f(X)$. \rd{} \tcr{(ref def)} implies that there exists $x \in X$ such that $f(x) = y$. Since $y \in Y$ is arbitrary, we have that for each $y \in Y$, there exists $x \in X$ such that $f(x) = y$. Hence $f$ is surjective.
			\end{itemize}
		\end{enumerate}
	\end{proof}
	
	
	
	
	
	
	
	
	
	
	
	
	
	
	
	
	
	
	
	
	
	
	
	
	
	
	
	
	
	
	
	
	
	
	
	
	
	
	
	
	
	
	
	
	
	
	
	
	
	
	
	
	
	
	\subsection{Axiom of Choice}
	
	\begin{defn}
		Let $A, B$ be sets, $F:A \rightarrow \MP(B)$ and $f: A \rightarrow B$. Then $f$ is said to be a \tbf{selection of $F$} if for each $a \in A$, $f(a) \in F(a)$.
	\end{defn}
	
	\begin{ax} \tbf{Axiom of Choice}: \\
		Let $A, B$ be sets, $F:A \rightarrow \MP(B)$. If for each $a \in A$, $F(a) \neq \varnothing$, then there exists $f:A \rightarrow B$ such that $f$ is a selection of $F$.  
	\end{ax}
	
	
	
	
	
	
	
	
	
	
	
	
	
	
	
	
	
	
	
	
	
	
	
	
	
	
	
	
	
	
	
	
	
	
	
	
	
	
	
	
	
	
	
	
	
	
	
	
	
	\subsection{Composition}
	
	\begin{defn}
		Let $X,Y,Z$ be sets, $f:X \rightarrow Y$ and $g:Y \rightarrow Z$. We define the \tbf{composition of $g$ with $f$}, denoted $g \circ f: X \rightarrow Y$, by 
		$$g \circ f(x) \defeq g(f(x)) \quad (x \in X)$$
	\end{defn}
	
	\begin{ex}
		Let $X,Y,Z$ be sets, $f:X \rightarrow Y$ and $g:Y \rightarrow Z$. Then $g \circ f: X \rightarrow Y$.
	\end{ex}
	
	\begin{proof}\
		\begin{enumerate}
			\item Let $x \in X$. Since $f:X \rightarrow Y$ and $x \in X$, $f(x) \in Y$. Define $y \in Y$ by $y \defeq f(x)$. Since $g:Y \rightarrow Z$ and $y \in Y$, $g(y) \in Z$. Therefore 
			\begin{align*}
				g \circ f(x)
				& = g(f(x)) \\
				& = g(y) \\
				& \in Z.
			\end{align*}
			Since $x \in X$ is arbitrary, we have that for each $x \in X$, $g \circ f(x) \in Z$. 
			\item Let $x_1, x_2 \in X$. Suppose that $x_1 = x_2$. Define $y_1, y_2 \in Y$ by $y_1 \defeq f(x_1)$ and $y_2 \defeq f(x_2)$. Since $f:X \rightarrow Y$ and $x_1 = x_2$, we have that 
			\begin{align*}
				y_1
				& = f(x_1) \\
				& = f(x_2) \\
				& = y_2.
			\end{align*} 
			Since $g:Y \rightarrow Z$ and $y_1 = y_2$, we have that 
			\begin{align*}
				g \circ f(x_1)
				& = g(f(x_1)) \\
				& = g(y_1) \\
				& = g(y_2) \\
				& = g(f(x_2)) \\
				& = g \circ f(x_2).
			\end{align*}
			Since $x_1, x_2 \in X$ with $x_1 = x_2$ are arbitrary, we have that for each $x_1, x_2 \in X$, $x_1 = x_2$ implies that $g \circ f(x_1) = g \circ f(x_2)$.
		\end{enumerate}
		Thus $g \circ f:X \rightarrow Y$. 
	\end{proof}
	
	\begin{ex}
		Let $X,Y,Z, W$ be sets, $f:X \rightarrow Y$, $g:Y \rightarrow Z$ and $h:Z \rightarrow W$. Then $(h \circ g) \circ f = h \circ (g \circ f)$.
	\end{ex}
	
	\begin{proof}
		For each $x \in X$,
		\begin{align*}
			(h \circ g) \circ f(x)
			& = (h \circ g) (f(x)) \\
			& = h(g(f(x))) \\
			& = h(g \circ f(x)) \\
			& = h \circ (g \circ f)(x).
		\end{align*}
		Thus $(h \circ g) \circ f = h \circ (g \circ f)$.
	\end{proof}
	
	\begin{defn}
		Let $X$ be a set. We define the \tbf{identity function on $X$}, denoted $\id_X:X \rightarrow X$ by $\id_X(x) \defeq x \quad (x \in X)$.
	\end{defn}
	
	\begin{ex}
		Let $X$ be a set. Then $\id_X:X \rightarrow X$. 
	\end{ex}
	
	\begin{proof}\
		\begin{enumerate}
			\item By definition, for each $x \in X$, 
			\begin{align*}
				\id_X(x)
				& = x \\
				& \in X.
			\end{align*}
			\item By definition, for each $x_1, x_2 \in X$, if $x_2 = x_2$, then 
			\begin{align*}
				\id_X(x_1)
				& = x_1 \\
				& = x_2 \\
				& = \id_X(x_2).
			\end{align*}
		\end{enumerate}
		Thus $\id_X:X \rightarrow X$. 
	\end{proof}
	
	\begin{ex}
		Let $X,Y$ be sets and $f:X \rightarrow Y$. Then 
		\begin{enumerate}
			\item $f \circ \id_X = f$, 
			\item $\id_Y \circ f = f$.
		\end{enumerate}
	\end{ex}
	
	\begin{proof}\
		\begin{enumerate}
			\item For each $x \in X$,
			\begin{align*}
				f \circ \id_X(x)
				& = f(\id_X(x)) \\
				& = f(x).
			\end{align*}
			Thus $f \circ \id_X = f$.
			\item For each $y \in Y$,
			\begin{align*}
				\id_Y \circ f(x)
				& = \id_Y(f(x)) \\
				& = f(x).
			\end{align*}
			Thus $\id_Y \circ f = f$.
		\end{enumerate}
	\end{proof}
	
	\begin{defn}
		Let $X,Y$ be sets and $f:X \rightarrow Y$. 
		\begin{itemize}
			\item Let $g:Y \rightarrow X$. Then $g$ is said to be a(n)
			\begin{itemize}
				\item \tbf{left inverse of $f$} if $g \circ f = \id_X$,
				\item \tbf{right inverse of $f$} if $f \circ g = \id_Y$,
				\item \tbf{inverse of $f$} if $g$ is a left inverse of $f$ and $g$ is a right inverse of $f$.
			\end{itemize}
			\item Then $f$ is said to 
			\begin{itemize}
				\item have a(n)
				\begin{itemize}
					\item \tbf{left inverse} if there exists $g:Y \rightarrow X$ such that $g$ is a left inverse of $f$,
					\item \tbf{right inverse} if there exists $g:Y \rightarrow X$ such that $g$ is a right inverse of $f$.
					\item \tbf{inverse} if there exists $g:Y \rightarrow X$ such that $g$ is an inverse of $f$,
				\end{itemize}
				\item be 
				\begin{itemize}
					\item \tbf{left invertible} if $f$ has a left inverse,
					\item \tbf{right invertible} if $f$ has a right inverse,
					\item \tbf{invertible} if $f$ has an inverse.
				\end{itemize}
			\end{itemize}
		\end{itemize} 
	\end{defn}
	
	\begin{ex}
		\tcr{make ex about nonunique left and right inverses}
	\end{ex}
	
	\begin{proof}
		\tcr{finish!!!}
	\end{proof}
	
	\begin{ex}
		Let $X,Y$ be sets, $f:X \rightarrow Y$, $g,g':Y \rightarrow X$. If $g$ is an inverse of $f$ and $g'$ is an inverse of $f$, then $g = g'$.
	\end{ex}
	
	\begin{proof}
		Suppose that $g$ is an inverse of $f$ and $g'$ is an inverse of $f$. Then 
		\begin{align*}
			g'
			& = g' \circ \id_Y \\
			& = g' \circ (f \circ g) \\
			& = (g' \circ f) \circ g \\
			& = \id_X \circ g \\
			& = g.
		\end{align*}
	\end{proof}
	
	\begin{defn}
		Let $X,Y$ be sets and $f:X \rightarrow Y$. Suppose that $f$ is invertible. We define the \tbf{inverse of $f$}, denoted $f^{-1}:Y \rightarrow X$, to be the unique $g:Y \rightarrow X$ such that $g$ is an inverse of $f$. 
	\end{defn}
	
	
	
	
	
	
	
	
	
	
	
	
	
	
	
	
	
	
	
	
	
	
	
	
	
	
	
	
	
	
	
	
	
	
	
	
	
	
	
	
	
	
	
	
	
	
	
	
	
	
	
	
	
	\subsubsection{Composition and X-jectivity}
	
	\begin{ex}
		Let $X$ be a set. Then
		\begin{enumerate}
			\item $\id_X$ is injective,
			\item $\id_X$ is surjective,
			\item $\id_X$ is bijective.
		\end{enumerate}
	\end{ex}
	
	\begin{proof}\
		\begin{enumerate}
			\item For each $x_1, x_2 \in X$, $\id_X(x_1) = \id_X(x_2)$ implies that 
			\begin{align*}
				x_1
				& = \id_X(x_1) \\
				& = \id_X(x_2) \\
				& = x_2. 
			\end{align*}
			Thus $\id_X$ is injective.
			\item Let $y \in X$. Define $x \in X$ by $x \defeq y$. Then 
			\begin{align*}
				\id_X(x)
				& = x \\
				& = y.
			\end{align*}
			Since $y \in X$ is arbitrary, we have that for each $y \in X$, there exists $x \in X$ such that $\id_X(x) = y$. Hence $\id_X$ is surjective.
			\item Since $\id_X$ is injective and $\id_X$ is surjective, we have that $\id_X$ is bijective.
		\end{enumerate}
	\end{proof}
	
	\begin{ex}
		Let $X,Y,Z$ be sets, $f:X \rightarrow Y$ and $g:Y \rightarrow Z$. 
		\begin{enumerate}
			\item If $g$ is injective and $f$ is injective, then $g \circ f$ is injective.
			\item If $g$ is surjective and $f$ is surjective, then $g \circ f$ is surjective.
			\item If $g$ is bijective and $f$ is bijective, then $g \circ f$ is bijective.
		\end{enumerate}
	\end{ex}
	
	\begin{proof}\
		\begin{enumerate}
			\item Suppose that $g$ is injective and $f$ is injective. Let $x_1, x_2 \in X$. Suppose that $g \circ f(x_1) = g \circ f(x_2)$. Define $y_1,y_2 \in Y$ by $y_1 \defeq f(x_1)$ and $y_2 \defeq f(x_2)$. Then 
			\begin{align*}
				g(y_1)
				& = g(f(x_1)) \\
				& = g \circ f(x_1) \\
				& = g \circ f(x_2) \\
				& = g(f(x_2)) \\
				& = g(y_2).
			\end{align*}
			Since $g$ is injective, we have that $y_1 = y_2$. Therefore
			\begin{align*}
				f(x_1)
				& = y_1 \\
				& = y_2 \\
				& = f(x_2).
			\end{align*}
			Since $f$ is injective, we have that $x_1 = x_2$. Since $x_1, x_2 \in X$ with $g \circ f(x_1) = g \circ f(x_2)$ are arbitrary, we have that for each $x_1, x_2 \in X$, $g \circ f(x_1) = g \circ f(x_2)$ implies that $x_1 = x_2$. Thus $g \circ f$ is injective.  
			\item Suppose that $g$ is surjective and $f$ is surjective. Let $z \in Z$. Since $g$ is surjective, there exists $y \in Y$ such that $g(y) = z$. Since $f$ is surjective, there exists $x \in X$ such that $f(x) = y$. Then 
			\begin{align*}
				g \circ f(x)
				& = g(f(x)) \\
				& = g(y) \\
				& = z.
			\end{align*}
			Since $z \in Z$ is arbitrary, we have that for each $z \in Z$, there exists $x \in X$ such that $g \circ f(x) = z$. Hence $g \circ f$ is surjective.
			\item 
			Suppose that $g$ is bijective and $f$ is bijective. Then $g$ is injective, $g$ is surjective, $f$ is injective and $f$ is surjective. \tcr{the prev parts (ref parts)} implies that $g \circ f$ is injective and $g \circ f$ is surjective. Thus $g \circ f$ is bijective.
		\end{enumerate}
	\end{proof}
	
	\begin{ex}
		Let $X,Y,Z$ be sets, $f:X \rightarrow Y$ and $g:Y \rightarrow Z$. Then
		\begin{enumerate}
			\item $g$ is injective implies that $g \circ f$ is injective iff $f$ is injective.
			\item $f$ is surjective implies that $g \circ f$ is surjective iff $g$ is surjective.
		\end{enumerate}
	\end{ex}
	
	\begin{proof}\
		\begin{enumerate}
			\item Suppose that $g$ is injective. 
			\begin{itemize}
				\item $(\implies)$: \\
				Suppose that $g \circ f$ is injective. Let $x_1, x_2 \in X$. Suppose that $f(x_1) = f(x_2)$. Define $y_1, y_2 \in Y$ by $y_1 \defeq f(x_1)$ and $y_2 \defeq f(x_2)$. By assumption,  
				\begin{align*}
					y_1
					& = f(x_1) \\
					& = f(x_2) \\
					& = y_2.
				\end{align*}
				Since $g:Y \rightarrow Z$, we have that $g(y_1) = g(y_2)$. Therefore
				\begin{align*}
					g \circ f(x_1) 
					& = g(f(x_1)) \\
					& = g(y_1) \\
					& = g(y_2) \\
					& = g(f(x_2)) \\
					& = g \circ f(x_2).
				\end{align*}
				Since $g \circ f$ is injective, we have that $x_1 = x_2$. Since $x_1, x_2 \in X$ with $f(x_1) = f(x_2)$ are arbitrary, we have that for each $x_1, x_2 \in X$, $f(x_1) = f(x_2)$ implies that $x_1 = x_2$. Thus $f$ is injective.
				\item $(\impliedby)$: \\
				Since $f$ is injective, \rex{} \tcr{(ref prev ex)} implies that if $f$ is injective, then $g \circ f$ is injective. 
			\end{itemize}
			\item Suppose that $f$ is surjective. 
			\begin{itemize}
				\item $(\implies)$: \\
				Suppose that $g \circ f$ is surjective. Let $z \in Z$. Since $g \circ f$ is surjective, there exists $x \in X$ such that $g \circ f(x) = z$. Define $y \in Y$ by $y \defeq f(x)$. Then 
				\begin{align*}
					g(y)
					& = g(f(x)) \\
					& = g \circ f(x) \\
					& = z. 
				\end{align*}
				Since $z \in Z$ is arbitrary, we have that for each $z \in Z$, there exists $y \in Y$ such that $g(y) = z$. Thus $g$ is surjective. 
				\item $(\impliedby)$: \\
				Since $f$ is surjective, \rex{} \tcr{(ref prev ex)} implies that if $g$ is surjective, then $g \circ f$ is surjective. 
			\end{itemize}
		\end{enumerate}
	\end{proof}
	
	\begin{ex}
		Let $X,Y$ be sets and $f:X \rightarrow Y$. Then 
		\begin{enumerate}
			\item $f$ is left invertible iff $f$ is injective.
			\item $f$ is right invertible iff $f$ is surective.
			\item $f$ is invertible iff $f$ is a bijection.
		\end{enumerate}
	\end{ex}
	
	\begin{proof}\
		\begin{enumerate}
			\item 
			\begin{itemize}
				\item $(\implies)$: \\
				Supose that $f$ is left invertible. Then $f$ has a left inverse. Thus there exists $g:Y \rightarrow X$ such that $g \circ f = \id_X$. Therefore for each $x_1, x_2 \in X$, $f(x_1) = f(x_2)$ implies that 
				\begin{align*}
					x_1
					& = \id_X(x_1) \\
					& = g \circ f(x_1) \\
					& = g(f(x_1)) \\
					& = g(f(x_2)) \\
					& = g \circ f(x_2) \\
					& = \id_X(x_2) \\
					x_2.
				\end{align*}
				Hence $f$ is injective.
				\item $(\impliedby)$: \\
				\begin{itemize}
					\item Suppose that $f$ is injective. \rex{} \tcr{(ref ex)} implies that for each $y \in f(X)$, there exists $x_y \in X$ such that $f^{-1}(\{y\}) = \{x_y\}$. Define $F:Y \rightarrow \MP(X)$ by $F(y) \defeq f^{-1}(\{y\})$. Since for each $y \in Y$, $F(y) \neq \varnothing$, the axiom of choice \tcr{(ref axiom)} implies that there exists $g:Y \rightarrow X$ such that for each $y \in Y$, $g(y) \in F(y)$. \rex{} \tcr{(ref ex)} implies that for each $y \in Y$, $g(y) = x_y$. 
					\item Let $x \in X$. Define $y \in Y$ by $y \defeq f(x)$. By definition of $g$, 
					\begin{align*}
						x_y
						& = g(y) \\
						& \in f^{-1}(\{y\}). 
					\end{align*}
					Thus 
					\begin{align*}
						f(x_y) 
						& = y \\
						& = f(x). 
					\end{align*}
					Since $f$ is injective, we have that 
					\begin{align*}
						x_{f(x)}
						& = x_y \\
						& = x.
					\end{align*}
					Since $x \in X$ is arbitrary, we have that for each $x \in X$, $x_{f(x)} = x$.
					\item Then for each $x \in X$, 
					\begin{align*}
						g \circ f(x)
						& = g(f(x)) \\
						& = x_{f(x)} \\
						& = x \\
						& = \id_X(x). 
					\end{align*} 
					Hence $g \circ f = \id_X$. Thus $f$ has a left inverse and therefore $f$ is left invertible. 
				\end{itemize}
			\end{itemize}
			\item 
			\begin{itemize}
				\item $(\implies)$: \\
				Supose that $f$ is right invertible. Then $f$ has a right inverse. Thus there exists $g:Y \rightarrow X$ such that $f \circ g = \id_Y$. 
				
				Hence $f$ is surjective.
				\item $(\impliedby)$: \\
				\begin{itemize}
					\item Suppose that $f$ is surjective. 
				\end{itemize}
			\end{itemize}
			\item 
		\end{enumerate}
	\end{proof}
	
	\begin{ex}
		Let $X,Y$ be sets and $f:X \rightarrow Y$. Then $f$ is invertible iff $f$ is a bijection.
	\end{ex}
	
	\begin{proof}
		\tcr{maybe make an exercise characterizing injectivity and surjectivity in terms of left and right inverses (basically the proof below) and then just reference that exercise here, need axiom of choice and previous exercise about product of singletons being singleton}
		\begin{itemize}
			\item $(\implies)$: \\
			Suppose that $f$ is invertible. 
			\begin{itemize}
				\item For each $x_1, x_2 \in X$, $f(x_1) = f(x_2)$ implies that 
				\begin{align*}
					x_1
					& = \id_X(x_1) \\
					& = f^{-1} \circ f(x_1) \\
					& = f^{-1}(f(x_1)) \\
					& = f^{-1}(f(x_2)) \\
					& = f^{-1} \circ f(x_2) \\
					& = \id_X(x_2) \\
					x_2.
				\end{align*}
				Thus $f$ is injective.
				\item Let $y \in Y$. Define $x \in X$ by $x \defeq f^{-1}(y)$. Then 
				\begin{align*}
					f(x)
					& = f(f^{-1}(y)) \\
					& = f \circ f^{-1}(y) \\
					& = \id_Y(y) \\
					& = y.
				\end{align*}
				Since $y \in Y$ is arbitrary, we have that for each $y \in Y$, there exists $x \in X$ such that $f(x) = y$. Hence $f$ is surjective.
			\end{itemize}
			\item $(\impliedby)$: \\
			Suppose that $f$ is a bijection. Then $f$ is injective and surjective. Since $f$ is surjective, $f(X) = Y$. Therefore for each $y \in Y$, $f^{-1}(\{y\}) \neq \varnothing$ \tcr{(define $f^{-1}(A)$ and $f(A)$, show $f(X) = Y$ iff $f$ is surjective and ref exercises)}. Since $f$ is injective and $f(X) = Y$, for each $y \in f(X)$, there exists $x_y \in X$ such that $f^{-1}(\{y\}) = \{x_y\}$. The axiom of choice \tcr{(actually an exercise I proved somewhere about products of singleton sets being a singelton set)} implies that there exists a unique $g:Y \rightarrow X$ such that for each $y \in Y$, $g(y) = x_y$. Then for each $y \in Y$, 
			\begin{align*}
				f \circ g(y) 
				& = f(g(y)) \\
				& = f(x_y) \\
				& = y \\
				& = \id_Y(y).
			\end{align*}
			Thus $f \circ g = \id_Y$. Let $x \in X$. Since $f$ is injective,
			\begin{align*}
				\{x\}
				& = f^{-1}(\{f(x)\}) \\
				& = \{g(f(x))\} \quad \text{(by definition of $g$)}.
			\end{align*}
			Therefore
			\begin{align*}
				g \circ f(x)
				& = g(f(x)) \\
				& = x \\
				& = \id_X(x).
			\end{align*}
			Since  $x \in X$ is arbitrary, we have that for each $x \in X$, $g \circ f(x) = \id_X(x)$. Hence $g \circ f = \id_X$.
		\end{itemize}
		Thus there exists $g:Y \rightarrow X$ such that $g \circ f = \id_X$ and $f \circ g = \id_Y$. Hence $f$ is invertible.  
	\end{proof}
	
	\begin{defn}
		Let $X, \ol{X}, Y$ be sets, $\iota:X \rightarrow \ol{X}$, $f:X \rightarrow Y$ and $\bar{f}: \ol{X} \rightarrow Y$. Then $\bar{f}$ is said to be an \tbf{extension of $f$ along $\iota$} if $\bar{f} \circ \iota = f$, i.e. the following diagram commutes:  
		\[ 
		\begin{tikzcd}
			X  \arrow[d, "\iota"'] \arrow[r, "f"] & Y \\
			\ol{X} \arrow[ur, "\bar{f}"'] 
		\end{tikzcd}
		\] 
	\end{defn}
	
	\begin{defn}
		Let $X, \ol{Y}, Y$ be sets, $\pi:\ol{Y} \rightarrow Y$, $f:X \rightarrow Y$ and $\bar{f}: X \rightarrow \ol{Y}$. Then $\bar{f}$ is said to be an \tbf{lift of $f$ through $\pi$} if $\pi \circ \bar{f} = f$, i.e. the following diagram commutes:  
		\[ 
		\begin{tikzcd}
			                                       & \ol{Y}  \arrow[d, "\pi"] \\
			X \arrow[r, "f"'] \arrow[ur, "\bar{f}"] & Y \\
		\end{tikzcd}
		\] 
	\end{defn}
	
	\tcr{Say some more about these and connect them back to previous material, eg retraction/sections being lifts/extensions ivolving the identity}
	
	\begin{defn}
		Let $X, \ol{X}, Y$ be sets, $\iota:X \rightarrow \ol{X}$, $f:X \rightarrow Y$ and $\bar{f}: \ol{X} \rightarrow Y$. We define the 
		\begin{itemize}
			\item \tbf{left factors of $f$ by $\iota$}, denoted $\LFactors(f, \iota) \subset Y^{\ol{X}}$, by $\LFactors(f, \iota) \defeq \{\bar{f} \in Y^{\ol{X}}: \text{$\bar{f}$ is an extension of $f$ along $\iota$}\}$.
			\item If there exists $\bar{f} \in Y^{\ol{X}}$ such that $\LFactors(f, \iota) = \{\bar{f}\}$, we define the \tbf{left factor of $f$ by $\iota$}, denoted $\LFactor(f, \iota) \in Y^{\ol{X}}$, by $\LFactor(f, \iota) \defeq \bar{f}$. 
		\end{itemize}
	\end{defn}
	
	\begin{defn}
		Let $X, \ol{X}, Y$ be sets, $\pi:\ol{Y} \rightarrow Y$, $f:X \rightarrow Y$ and $\bar{f}: X \rightarrow \ol{Y}$. 
		\begin{itemize}
			\item We define the \tbf{right factors of $f$ by $\pi$}, denoted $\LFactors(f, \pi) \subset Y^{\ol{X}}$, by $\Factors(f, \pi) \defeq \{\bar{f} \in {\ol{Y}}^X: \text{$\bar{f}$ is a lift of $f$ through $\pi$}\}$.
			\item If there exists $\bar{f} \in {\ol{Y}}^X$ such that $\RFactors(f, \pi) = \{\bar{f}\}$, we define the \tbf{right factor of $f$ by $\pi$}, denoted $\RFactor(f, \pi) \in {\ol{Y}}^X$, by $\RFactor(f, \pi) \defeq \bar{f}$. 
		\end{itemize}
	\end{defn}
	
	
	
	
	
	
	
	
	
	
	
	
	
	
	
	
	
	
	
	
	
	
	
	
	
	
	\subsection{Schroder-Bernstein Theorem:}
	
	\tcr{Motivate Schroder-Bernstein theorem with Hilbert's hotel and reference blargoner's youtube video \href{https://www.youtube.com/watch?v=IkoKttTDuxE}{https://www.youtube.com/watch?v=IkoKttTDuxE}}
	
	\tcr{Write some stuff about hilbert's hotel}
	
	\begin{questn}
		Suppose we have a hotel with infintely many rooms, each of which is current occupied by a unique guest. Suppose infintely many more prospective guests arrive in search of a room. 
		
		Now assume that 
		\begin{itemize}
			\item there is a room-assignment procedure which will assign a room to each of the guests (prospective or current) but will not assign more than one guest to the same room,
			\item there is a occupant-identification procedure which for each room, allows us to identify the current occupant.
		\end{itemize}
		
		We note that the room assingment procedure may not necessarily fill all the rooms and the occupant identification procedure may not necessarily provide the names of all guests (e.g. the prospective guests before they are assigned a room).
		
		Is it possible to assign a room to each guest (current or prospective) using the room-assignment and occupant-identification procedures in such a way so that each of the rooms in the hotel are occupied 
	\end{questn}
	
	\tcr{draw picture}
	
	\begin{soln}
		We will denote the set of guests (current or prospective) by $Y$ and the set of rooms by $X$. We denote the room-assignment procedure by $g:Y \rightarrow X$ and we denote the occupant-identification procedure by $f:X \rightarrow Y$. 
		
		We may begin by 
		\begin{enumerate}
			\item assigning each of the prospective guests to a room using the room-assingment procedure, that is we will send each $y \in f(X)^c$ to $g(y) \in g[f(X)^c]$.
			\item Since the rooms in $g[f(X)^c]$ are already currently occupied, for each room in $g[f(X)^c]$, we may identify its occupant using the occupant-identification procedure and temporarily remove them from their room, that is send each $x \in g[f(X)^c]$ to $f(x) \in (f \circ g)[f(X)^c]$.
		\end{enumerate}
		We may then inductively repeat the previous steps so that after the $k$-th iteration of the previous steps, we have just assigned each guest $y \in (f \circ g)^{k-1}[f(X)^c]$ to room $g(y) \in g \circ (f \circ g)^{k-1}[f(X)^c]$ and temporarily removed each guest $f(x) \in (f \circ g)^k[f(X)^c]$. 
	\end{soln}
	
	\tcr{rework intuition, move intuition about size to exercises about injectivity and surjectivity, draw pictures, add picture and intuition about same size-ness for bijections, draw strings, etc}
	
	Let $X$ and $Y$ be sets. Suppose that there exists $f:X \rightarrow Y$ such that $f$ is injective.  Intuitively, we are able to consider $Y$ as being ``at least as large as" $X$ because for each $x \in X$, we may tie a string from $x$ to a unique $f(x) \in Y$, so that there are ``at least as many" elements of $Y$ as there are elements in $X$. In this analogy, it is not generally the case that each $y \in Y$ will have a string attached to it unless $f$ is a bijection. However, what if we find ourselves in the situation in which there also exists $g:Y \rightarrow X$ such that $g$ is an injection. Then intuitively, $X$ ``at least as large as" $Y$ and we would then want it to be the case that $X$ and $Y$ are ``the same size". Taking our analogy further, we would like it to be the case that there is a bijection $h:X \rightarrow Y$. The previous question about hotels provides the mechanism for showing this to be the case.
	
	\begin{ex} \tbf{Schroder-Bernstein Theorem:} \\
		Let $X,Y$ be sets. If there exist $f:X \rightarrow Y$, $g:Y \rightarrow X$ such that $f$ and $g$ are injections, then there exists $\phi:X \rightarrow Y$ such that $\phi$ is a bijection.
	\end{ex}
	
	\begin{proof}
		Define $A_0 \defeq f(X)^c$ and $A \defeq \bigcup\limits_{k=0} (f \circ g)^{\circ k} (A_0)$. \tcr{Show} $A^c \subset f(X)$ and $f(A^c) = g(A)^c$ and \tcr{show} $g|^{g(A)}_A$ and $(f|^{A^c}_{g(A)^c})^{-1}$ are invertible. Define $h:Y \rightarrow X$ by 
		\[
		h(y) \defeq 
		\begin{cases}
		g(y), & y \in A \\
		(f|^{f(X)})^{-1}(y), & y \in A^c	
		\end{cases}
		\]
		i.e. $h = g|^{g(A)}_A \oplus (f|^{A^c}_{g(A)^c})^{-1}$\\
		\tcr{there may be an exercise showing that inverse of coprod is coprod of inverses}
	\end{proof}
	
	
	
	
	
	
	
	
	
	
	
	
	
	
	
	
	
	
	
	
	
	
	
	
	
	
	
	
	
	
	
	
	
	
	
	
	
	
	
	
	
	
	
	
	
	
	
	
	
	
	
	
	
	
	
	
	
	
	
	
	
	
	
	
	
	\section{Equivalence Relations}
	
	\begin{note}
		Let $X$ be a set and $\nu \in \MR^2(X)$. When the context is clear, we may write $a \nu b$ in place of $(a,b) \in \nu$.
	\end{note}
	
	\begin{defn}
		Let $X$ be a set. Suppose that $X \neq \varnothing$.
		\begin{itemize}
			\item Let $\th \in \MR^2(X)$. Then $\th$ is said to be 
			\begin{itemize}
				\item \tbf{reflexive} if for each $a \in X$, $(a,a) \in \th$
				\item \tbf{symmetric} if for each $a,b \in X$, $(a,b) \in \th$ implies that $(b,a) \in \th $
				\item \tbf{transitive} if for each $a,b,c \in X$, $(a,b) \in \th$ and $(b,c) \in \th$ implies that $(a,c) \in \th $
				\item an \tbf{equivalence relation on $X$} if $\th$ is reflexive, symmetric and transitive
			\end{itemize}
			\item We define $\Equi(X) \defeq \{\th \in \MR^2(X): \text{$\th$ is an equivalence relation on $X$}\}$
		\end{itemize}
	\end{defn}
	
	\tcr{Make some exercise about equiv relations}
	
	\begin{defn}
		Let $X$ be a set and $\th \in \Equi(X)$.
		\begin{itemize}
			\item Let $x \in X$. We define the \tbf{equivalence class of $x$ under $\th$}, denoted $[x]_{\th} \subset X$, by $[x]_{\th} \defeq \{y \in X: (x,y) \in \th\}$.  
			\item We define the \tbf{quotient set of $X$ by $\th$}, denoted $X/\th \subset \MP(X)$, by $X/\th \defeq \{[x]_{\th}: x \in X\}$. 
		\end{itemize}
	\end{defn}
	
	\begin{note}
		When the context is clear, we may write $[x]$ or $\bar{x}$ in place of $[x]_{\th}$.
	\end{note}
	
	\begin{ex}
		Let $X$ be a set and $\th \in \Equi(X)$. 
		\begin{enumerate}
			\item For each $x \in X$, $[x]_{\th} \neq \varnothing$.
			\item For each $x,y \in X$, $(x,y) \in \th$ iff $[x]_{\th} = [y]_{\th}$.
			\item For each $x,y \in X$, $(x,y) \not \in \th$ iff $[x]_{\th} \cap [y]_{\th} = \varnothing$.
		\end{enumerate}
	\end{ex}
	
	\begin{proof}\
		\begin{enumerate}
			\item Let $x \in X$. Since $\th$ is reflexive, $(x,x) \in \th$. Thus $x \in [x]_{\th}$. Hence $[x]_{\th} \neq \varnothing$. Since $x \in X$ is arbitrary, we have that for each $x \in X$, $[x]_{\th} \neq \varnothing$.
			\item Let $x,y \in X$.
			\begin{itemize}
				\item $(\implies)$: \\
				Suppose that $(x,y) \in \th$. Then for each $z \in X$, 
				\begin{align*}
					z \in [x]_{\th}
					& \iff (x,z) \in \th \\
					& \iff (y,z) \in \th \quad \text{(since $(x,y) \in \th$ and $\th$ is transitive)} \\
					& \iff z \in [x]_{\th}.
				\end{align*}
				Hence $[x]_{\th}  = [y]_{\th}$. 
				\item $(\impliedby)$: \\
				We have that 
				\begin{align*}
					[x]_{\th} = [y]_{\th}
					& \implies y \in [x]_{\th} \quad \text{(since $y \in [y]_{\th}$)} \\
					& \implies (x,y) \in \th. 
				\end{align*} 
			\end{itemize} 
			\item Let $x,y \in X$.
			\begin{itemize}
				\item $(\implies)$: \\
				Suppose that $(x,y) \not \in \th$. For the sake of contradiction, suppose that $[x]_{\th} \cap [y]_{\th} \neq \varnothing$. Then there exists $z \in X$, such that $z \in [x]_{\th} \cap [y]_{\th}$. We note that
				\begin{align*}
					z \in [x]_{\th} \cap [y]_{\th}
					& \iff \text{$z \in [x]_{\th}$ and $z \in [y]_{\th}$} \\
					& \iff \text{$(x,z) \in \th$ and $(y,z) \in \th$} \\
					& \iff \text{$(x,z) \in \th$ and $(z,y) \in \th$} \quad \text{(since $\th$ is symmetric)} \\
					& \implies (x,y) \in \th \quad \text{(since $\th$ is transitive)}.
				\end{align*}
				Hence $(x,y) \in \th$, which is a contradiction. Thus $[x]_{\th} \cap [y]_{\th} = \varnothing$.
				\item $(\impliedby)$: \\
				Suppose that $[x]_{\th} \cap [y]_{\th} = \varnothing$. Since $y \in [y]_{\th}$, we have that $y \not \in [x]_{\th}$. Thus $(x,y) \not \in \th$. 
			\end{itemize} 
		\end{enumerate}
	\end{proof}
	
	\begin{defn}
		Let $X$ be a set.
		\begin{itemize}
			\item Let $\MP \subset \MP(X)$. Then $\MP$ is called a \tbf{partition of $X$} if 
			\begin{enumerate}
				\item for each $E \in \MP$, $E \neq \varnothing$, 
				\item for each $E_1, E_2 \in \MP$, $E_1 \neq E_2$ implies that $E_1 \cap E_2 = \varnothing$,
				\item $\bigcup\limits_{E \in \MP} E = X$.
			\end{enumerate}
			\item We define $\Part(X) \defeq \{\MP \subset \MP(X): \text{$\MP$ is a partition of $X$}\}$.
		\end{itemize}
	\end{defn}
	
	\begin{ex}
		Let $X$ be a set and $\th \in \Equi(X)$. Then $X/\th \in \Part(X)$. 
	\end{ex}
	
	\begin{proof}\
		\begin{enumerate}
			\item Let $E \in X/\th$. Then there exists $x \in X$ such that $E = [x]_{\th}$. Then 
			\begin{align*}
				x
				& \in [x]_{\th} \\
				& = E
			\end{align*}
			and $E \neq \varnothing$. Since $E \in X/\th$ is arbitrary, we have that for each $E \in X/\th$, $E \neq \varnothing$.
			\item Let $E_1, E_2 \in X/\th$. Suppose that $E_1 \neq E_2$. Since $E_1, E_2 \in X/\th$, there exist $x_1, x_2 \in X$ such that $E_1 = [x_1]_{\th}$ and $E_2 = [x_2]_{\th}$. Since 
			\begin{align*}
				E_1 \neq E_2
				& \iff [x_1]_{\th} \neq [x_2]_{\th} \\
				& \iff (x_1, x_2) \not \in \th \quad \text{(by \tcr{prev ex})} \\
				& \iff [x_1]_{\th} \cap [x_2]_{\th} = \varnothing \quad \text{(by \tcr{prev ex})} \\
				& \iff E_1 \cap E_2 = \varnothing, 
			\end{align*} 
			we have that $E_1 \cap E_2 = \varnothing$. 
			\item Let $x \in X$. Define $E \in X/\th$ by $E \defeq [x]_{\th}$. Since $\th$ is reflexive, $x \in [x]_{\th}$. Thus there exists $E \in X/\th$ such that $x \in E$. Thus $x \in \bigcup\limits_{E' \in X/\th} E'$. Since $x \in X$ is arbitrary, we have that for each $x \in X$, $x \in \bigcup\limits_{E' \in X/\th} E'$. Hence $X \subset \bigcup\limits_{E' \in X/\th} E'$. Since $\bigcup\limits_{E' \in X/\th} E' \subset X$, we have that $X =  \bigcup\limits_{E' \in X/\th} E'$. 
		\end{enumerate}
		Thus $X/\th \in \Part(X)$.  
	\end{proof}
	
	\begin{defn}
		Let $X$ be a set and $\MP \in \Part(X)$. Define $\th_{\MP} \in \MR^2(X)$ by $\th_{\MP} \defeq \{(x,y) \in X^2: \text{there exists $E \in \MP$ such that $x,y \in E$}\}$.
	\end{defn}
	
	\begin{ex}
		Let $X$ be a set and $\MP \in \Part(X)$. Then 
		\begin{enumerate}
			\item $\th_{\MP} \in \Equi(X)$
			\item for each $x \in X$ and $F \in \MP$, $[x]_{\th_{\MP}} = F$ iff $x \in F$. 
			\item $X/\th_{\MP} = \MP$.
		\end{enumerate} 
	\end{ex}
	
	\begin{proof}\
		\begin{enumerate}
			\item 
			\begin{enumerate}
				\item Let $x \in X$. Since $\MP \in \Part(X)$, $X = \bigcup\limits_{E \in \MP} E$. Thus there exists $E \in \MP$ such that $x \in E$. Thus $(x,x) \in \th_{\MP}$. Since $x \in X$ is arbitrary, we have that for each $x \in X$, $(x,x) \in \th_{\MP}$. Thus $\th_{\MP}$ is reflexive. 
				\item We note that for each $x,y \in X$, 
				\begin{align*}
					(x,y) \in \th_{\MP}
					& \iff \text{there exists $E \in \MP$ such that $x,y \in E$} \\
					& \iff \text{there exists $E \in \MP$ such that $y,x \in E$} \\
					& \iff (y, x) \in \th_{\MP}. 
				\end{align*}
				Thus $\th_{\MP}$ is symmetric.
				\item Let $x,y,z \in X$. Suppose that $(x,y) \in \th_{\MP}$ and $(y,z) \in \th_{\MP}$. Then there exists $E, F \in \MP$ such that $x,y \in E$ and $y,z \in F$. Then $y \in E \cap F$. Since $E \cap F \neq \varnothing$, \tcr{by def} implies that $E = F$. Thus $x,z \in E$. Hence $(x,z) \in \th_{\MP}$. Since $x,y,z \in X$ with $(x,y) \in \th_{\MP}$ and $(y,z) \in \th_{\MP}$ are arbitrary, we have that for each $x,y,z \in X$, $(x,y) \in \th_{\MP}$ and $(y,z) \in \th_{\MP}$ implies that $(x,z) \in \th_{\MP}$. Hence $\th_{\MP}$ is transitive. 
			\end{enumerate}
			Thus $\th_{\MP} \in \Equi(X)$. 
			\item Let $x \in X$ and $F \in \MP$. 
			\begin{itemize}
				\item $[x]_{\th_{\MP}} = F$ iff $x \in F$. 
				\item 
			\end{itemize}
			\item Let $E \in X/\th_{\MP}$. Since $X/\th_{\MP} \in \Part(X)$, $E \neq \varnothing$. Thus there exists $x \in X$ such that $x \in E$. By \tcr{def of $X/\th_{\MP}$}, $E = [x]_{\th_{\MP}}$. Since $\MP \in \Part(X)$, $X = \bigcup\limits_{F \in \MP} F$. Thus there exists $F_x \in \MP$ such that $x \in F_x$. Since $\MP \in \Part(X)$, for each $F_1, F_2 \in \MP$, $x \in F_1 \cap F_2$ implies that $F_1 = F_2$. Thus $\{F \in \MP: x \in F\} = \{F_x\}$. By \tcr{def of $\th_{\MP}$},  
			\begin{align*}
				E
				& = [x]_{\th_{\MP}} \\
				& = \{y \in X: (x,y) \in \th_{\MP}\} \\
				& = \{y \in X: \text{there exists $F \in \MP$ such that $x,y \in F$}\} \\
				& = \{y \in X: y \in F_x\} \\
				& = F_x \\
				& \in \MP.
			\end{align*}
			Since 
			\tcr{FINISH!!!!}
			Let $y \in [x]_{\th_{\MP}}$.  $F \in \MP$. If $x,y \in F$, then $d$ $X = \bigcup\limits_{F}$
		\end{enumerate}
	\end{proof}
	
	\begin{ex}
		Let $X$ be a set. Then there is a $F: \Equi(X) \rightarrow \Part(X)$ such that $F$ is a bijection. \tcr{In Cat theory notes, show this is a natural iso between $\Equi$ and $\Part$ which are contravariant and for $f:X \rightarrow Y$, $\Part(f): \Part(Y) \rightarrow \Part(X)$ is given by $\Part(f)(P) \defeq f^*\MP$}
	\end{ex}
	
	\begin{proof}
		Define $F: \Equi(X) \rightarrow \Part(X)$ by $F(\th) \defeq X/\th$. 
		\begin{itemize}
			\item \tbf{Injectivity:} \\
			Let $\th_1, \th_2 \in \Equi(X)$. Suppose that $F(\th_1) = F(\th_2)$. Let $x,y \in X$. Then 
			\begin{align*}
				[x]_{\th_1},[y]_{\th_1} 
				& \in X/\th_1 \\
				& = F(\th_1) \\
				& = F(\th_2) \\
				& = X/\th_2. 
			\end{align*}
			Since $x \in [x]_{\th_1} \cap [x]_{\th_2}$, we have that $[x]_{\th_1} \cap [x]_{\th_2} \neq \varnothing$. \tcr{prev ex} implies that $X/\th_2 \in \Part(X)$. \tcr{By prev def}, $[x]_{\th_1} = [x]_{\th_2}$. Similarly, $[y]_{\th_1} = [y]_{\th_2}$. Since $x,y \in X$ are arbitrary, we have that for each $x,y \in X$ $[x]_{\th_1} = [x]_{\th_2}$ and $[y]_{\th_1} = [y]_{\th_2}$. Then for each $x,y \in X$,
			\begin{align*}
				(x,y) \in \th_1
				& \iff [x]_{\th_1} = [y]_{\th_1} \quad \text{(by prev ex)}\\
				& \iff [x]_{\th_2} = [y]_{\th_2} \quad \text{(from \tcr{above})} \\
				& \iff (x,y) \in \th_2.
			\end{align*}
			Thus $\th_1 = \th_2$. Since $\th_1, \th_2 \in \Equi(X)$ with $F(\th_1) = F(\th_2)$ are arbitrary, we have that for each $\th_1, \th_2 \in \Equi(X)$, $F(\th_1) = F(\th_2)$ implies that $\th_1 = \th_2$. Hence $F$ is injective.
			\item \tbf{Surjectivity:} \\
			Let $\MP \in \Part(X)$. \tcr{A prev ex} implies that $\th_{\MP} \in \Equi(X)$. Let $E \in X/\th_{\MP}$. Then  
		\end{itemize}      
	\end{proof}
	
	
	
	
	
	
	
	
	
	
	
	
	
	
	
	
	
	
	
	
	
	
	
	
	
	
	
	
	
	
	
	
	
	
	
	
	
	
	
	
	
	
	
	
	
	
	
	
	
	
	
	
	
	
	
	
	
	
	
	
	
	
	
	
	
	
	
	
	
	\chapter{Arithmetic}
	
	\subsection{Construction of Natural Numbers}
	
	
	
	\begin{defn} 
		We define the successor function, denoted $S: \bSet \rightarrow \bSet$, by $S(x) \defeq x \cup \{x\}$. 
	\end{defn}
	
	\begin{ax} \tbf{Axiom of Infinity}
		There exists $X$ such that $\varnothing \in X$ and for each $x \in X$, $S(x) \in X$. 
	\end{ax}
	
	\begin{defn} \tbf{Natural Numbers}
		We define $\N_0$ inductively by $0 \defeq \varnothing$ and $\N_0 \defeq \cl(\{0\}, S)$ \tcr{(define closure under maps in closure operator section of algebra notes)}. 
	\end{defn}
	
	\begin{note}
		Here we assign the symbols $1$ to $S(0)$, $2$ to $S(1)$, etc. 
	\end{note}
	
	\begin{defn}
		We define $+: \N_0 \rightarrow \N_0$ by $m+n \defeq S(m+n)$
	\end{defn}
	
	
	
	
	
	
	
	
	
	
	
	
	
	
	
	
	
	
	
	
	
	
	
	
	
	
	
	
	
	
	
	
	
	
	
	
	
	
	
	
	
	
	
	
	
	
	
	
	
	
	
	
	
	
	
	
	
	
	
	
	
	
	
	
	
	
	
	
	
	
	\part{Algebra}
	
	\chapter{Groups}
	
	\chapter{Rings}
	
	\chapter{Vector Spaces}
	
	
	
	
	
	
	
	
	
	
	
	
	
	
	
	
	
	
	
	
	
	
	
	
	
	
	
	
	
	
	
	
	
	
	
	
	
	
	
	
	
	
	
	
	
	
	
	
	
	
	
	
	
	
	
	
	
	
	
	
	
	
	
	
	
	
	
	\part{Analysis}
	
	\chapter{Real Numbers}
	
	\section{Construction and Properties of $\R$}
	
	\begin{defn}
		\tcr{least upper bound property}
	\end{defn}
	
	\begin{defn}
		Dedekind cuts
	\end{defn}
	
	\begin{ex} \tbf{The $\ep$-principle:} \\
		Let $x, y \in \R$. Then $x \leq y$ iff for each $\ep > 0$, $x \leq y = \ep$. 
	\end{ex}
	
	\begin{proof}\
		\begin{itemize}
			\item $(\implies)$: \\
			Suppose that $x \leq y$. Then for each $\ep >0$,
			\begin{align*}
				x
				& \leq y \\
				& = y + 0 \\
				& < y + \ep.
			\end{align*}
			\item $(\impliedby)$: \\
			Suppose that $x > y$. Define $\ep \in \R$ by $\ep \defeq (x - y)/2$. Then $\ep > 0$ and
			\begin{align*}
				y + \ep 
				& = y + \frac{x - y}{2} \\
				& = \frac{x + y}{2}.
			\end{align*} 
			Since $x > y$, \tcr{a prev ex about arithmetic mean} implies that
			\begin{align*}
				y + \ep 
				& = \frac{x + y}{2} \\
				& < x.
			\end{align*}
			Thus 
			$$(\text{$x > y$}) \implies (\text{$\exists \ep > 0$ such that $x > y + \ep$}).$$
			By contrapositive, we have that
			$$(\text{$\forall \ep > 0$, $x \leq y + \ep$}) \implies (\text{$x \leq y$}).$$
		\end{itemize}
	\end{proof}
	
	
	
	\section{Continuity}
	
	A notion of `closeness' constitutes a significant portion of the the foundation of analysis. As with other constructions in mathematics, the maps that respect or 'preserve' the notion we care about (e.g. closeness) are of great importance and utility when building our understanding of the objects and spaces under consideration (e.g. real numbers, minimum values of real-valued functions, etc). 
	
	In this section, we discuss maps $f: A \rightarrow B$ where $A,B \subset \R$, that preserve a notion of closeness on $\R$ in the sense that if points $(x_n)_{n \in \N} \subset A$ 'approach' a point $x_0 \in A$ as the index $n \in \N$ increases, then the images $(f(x_n))_{n \in \N} \subset B$ of $(x_n)_{n \in \N}$ will approach the image $f(x_0) \in B$ of $x_0$.  
	
	\section{Sequences}
	
	\section{Series}
	
	\begin{defn}
		Define $\sum\limits_{\al \in A} x_{\al}$ for $(x_{\al})_{\al \in A} \subset \R_+$ and $l^1(A)$.
	\end{defn}
	
	\begin{ex}
		\tcr{basic exercises about convergence of series}
	\end{ex}
	
	\section{Functions}
	
	
	
	
	
	
	
	
	
	
	
	
	
	
	
	
	
	
	
	
	
	
	
	
	
	
	
	
	
	
	
	
	
	
	
	
	
	
	
	
	
	
	
	
	
	
	\section{Limits of Functions}
	
	\begin{defn}
		Let $a,b \in \R$. Suppose that $a < b$. 
		\begin{itemize}
			\item We define $\Del_{(a,b)} \defeq \{(s,t) \in (a,b)^2: s < t\}$.
			\item Let $f \in \Bnd((a,b), \R)$. 
			\begin{itemize}
				\item For each $(s,t) \in \Del_{(a,b)}$, we define $M^f_{s,t}, m^f_{s,t} \in \R$ by $M^f_{s,t} \defeq \sup\limits_{x \in [s,t]} f(x)$ and $m^f_{s,t} \defeq \inf\limits_{x \in [s,t]} f(x)$
				\item We define $M^f, m^f: \Del_{(a,b)} \rightarrow \R$ by $M^f(s,t) \defeq M^f_{s,t}$ and $m^f(s,t) \defeq m^f_{s,t}$.
			\end{itemize}
		\end{itemize}
	\end{defn}
	
	\begin{ex} \lex{ex:analysis:real_nums:limits_funct:0010}
		Let $f \in \Bnd((a,b), \R)$ and $x_0 \in (a,b)$. If $f$ is continuous at $x_0$, then 
		\begin{enumerate}
			\item $\limt{x_0^+} M^f_{x_0, t} = f(x_0)$ 
			\item $\limt{x_0^+} m^f_{x_0, t} = f(x_0)$
		\end{enumerate}
	\end{ex}
	
	\begin{proof} 
		Suppose that $f$ is continuous at $x_0$. 
		\begin{enumerate}
			\item Let $\ep > 0$. Then there exists $\del > 0$ such that for each $x \in (a,b)$, $|x-x_0| < \del$ implies that $|f(x) - f(x_0)| < \ep/2$. Let $t \in (a, b)$. Suppose that $0 < t - x_0 < \del$. Since 
			\begin{align*}
				M^f_{x_0, t} - \frac{\ep}{2}
				& < M^f_{x_0, t} \\
				& = \sup\limits_{x \in [x_0, t]} f(x), 
			\end{align*}
			there exists $x' \in [x_0, t]$ such that $M^f_{x_0, t} - \ep/2 < f(x')$. Since 
			\begin{align*}
				|x' - x_0| 
				& = x' - x_0 \\
				& < t - x_0 \\
				& < \del,
			\end{align*}
			we have that $|f(x') - f(x_0)| < \ep/2$. Hence
			\begin{align*}
				|M^f_{x_0, t} - f(x_0)|
				& \leq |M^f_{x_0, t} - f(x')| + |f(x') - f(x_0)| \\
				& = (M^f_{x_0, t} - f(x')) + |f(x') - f(x_0)| \\
				& < \frac{\ep}{2} + \frac{\ep}{2} \\
				& = \ep.
			\end{align*}
			Since $t \in (a,b)$ with $0 < t - x_0 < \del$ is arbitrary, we have that for each $t \in (a,b)$, $0 < t - x_0 < \del$ implies that $|M^f_{x_0, t} - f(x_0)| < \ep$. Since $\ep > 0$ is arbitrary, we have that for each $\ep > 0$, there exists $\del > 0$ such that for each $t \in (a,b)$, $0 < t - x_0 < \del$ implies that $|M^f_{x_0, t} - f(x_0)| < \ep$. Thus $\limt{x_0^+} M^f_{x_0, t} = f(x_0)$.
			\item Similar to \tcr{prev part}.
		\end{enumerate}
	\end{proof}
	
	
	
	
	
	
	
	\chapter{Differentiation}
	
	\section{L'H\^opital's Rule}

	In \tcr{ref prev section on limits of sequences and functions and make/ref ex}, we considered sequences and functions like $(a_n)_{n \in \N} \subset \R$ and $f:(0,2) \rightarrow \R$ defined by 
	$$a_n \defeq  \frac{6n^2 + n +2}{2n^2 - \cos n}, \quad \quad f(x) \defeq \frac{x^2-4}{x-2}.$$ 
	Even though we were not able to immediately apply results like \tcr{ref ex about products/quotients/sums of limits} to $(a_n)_{n \in \N}$ and $f$, we were still able to deduce that $\limn a_n = 3$ and $\limneqx{2} f(x) = 4$ by reasoning about limiting behavior or exploiting simple algebraic structure and subsequently employing those results or using the definition of a limit. 
	
	There are many cases however, in which such an approach will not be practical. There are some situations where it may well be difficult to even intuit what the limiting behavior could be. For example, in \tcr{ref section on trig functions}, more effort was required in order to show that $\limneqx{0} x^{-1} \sin x = 1$. In this section, we will develop a tool that will enable us to handle a larger family of such cases. 
	
	To motivate the result, suppose that we are given $a,b \in \R$, $x_0 \in (a,b)$, $L \in \R$ and two functions $f,g:(a,b) \rightarrow \R$ which are differentiable and satisfy $|g'| > 0$ and $f'(x)/g'(x) \rightarrow L$ as $x \rightarrow x_0$. If $f(x) \rightarrow 0$ as $x \rightarrow x_0$ and $g(x) \rightarrow 0$ as $x \rightarrow x_0$, then intuitively, we have that for $x \approx x_0$ and $h \approx 0$, 
	\begin{align*}
		\frac{f(x + h)}{g(x + h)}
		& \approx \frac{f(x) + f'(x)h}{g(x) + g'(x)h} \\
		& \approx \frac{f'(x)h}{g'(x)h} \\
		& = \frac{f'(x)}{g'(x)} \\
		& \approx L.
	\end{align*}
	
	We formalize this intuition with the result below. 
	
	\begin{ex} \tbf{L'H\^opital's Rule}
		Let $a,b \in \RE$ and $f,g: (a,b) \rightarrow \R$. Suppose that $a < b$, $f,g$ are differenitable and for each $x \in (a,b)$, $g'(x) \neq 0$. Let $c \in [a,b]$. Suppose that $\limneqx{c} f'(x)/g'(x)$ exists in $\RE$. 
		\begin{enumerate}
			\item If $\limneqx{c} f(x) = 0$ and $\limneqx{c} g(x) = 0$, then 
			$$\limneqx{c} \frac{f(x)}{g(x)} = \limneqx{c} \frac{f'(x)}{g'(x)}$$
			\item If $\limneqx{c} |f(x)| = \infty$ and $\limneqx{c} |g(x)| = \infty$, then 
			$$\limneqx{c} \frac{f(x)}{g(x)} = \limneqx{c} \frac{f'(x)}{g'(x)}$$
		\end{enumerate}
	\end{ex}
	
	\begin{proof}\
		\begin{enumerate}
			\item Suppose that $\limneqx{c} f(x) = 0$ and $\limneqx{c} g(x) = 0$. 
			\item 
		\end{enumerate}
	\end{proof}
	
	\begin{note}
		\begin{itemize}
			\item Discuss why approximating $\frac{f(x + h)}{g(x + h)}$ with $\frac{f'(x)}{g'(x)}$ and getting a common denominator for the difference is somehow 'less good' approximation than what is doen in the proof, or even in the simple case of $x_0 \in (a,b)$ with $f(x_0) = g(x_0) = 0$, using a similar approximation of 
			\begin{align*}
				\frac{f(x_0 + h)}{g(x_0 + h)}
				& = \frac{f(x_0) + f'(x_0)h + o(|h|)}{g(x_0) + g'(x_0)h + o(|h|)} \\
				& = \frac{f'(x_0)h + o(|h|)}{g'(x_0)h + o(|h|)} \\
				& = \frac{|h|}{|h|}\frac{f'(x_0) + o(1)}{g'(x_0) + o(1)} \\
				& = \frac{f'(x_0) + o(1)}{g'(x_0) + o(1)} \\
				& \text{$\rightarrow \frac{f'(x_0) }{g'(x_0)}$ as $h \rightarrow 0$} 
			\end{align*}
			readily gives the result. The bad approximation ends with a $g(x+h)g'(x)$ in the denominator, among other problematic terms. Is there a way to quantify how or why the proof method
			\item \tcr{Maybe describe the controversy about how Lhopital didnt actally discover this, he was just wealthy and paid talented people and thanks to the continued PR, is now somewhat immortalized among math students}
		\end{itemize}
	\end{note}
	
	
	
	
	
	
	
	
	
	
	
	
	
	
	
	
	
	
	
	
	
	
	
	
	
	
	
	
	
	
	
	
	
	
	
	
	
	
	
	
	
	
	
	
	
	\chapter{Integration}
	
	\section{The Darboux Integral}
	
	\begin{defn} \ld{00000} 
		Let $a,b \in \R$. Suppose that $a<b$. Define $$\Bnd([a,b], \R) \defeq \{f:[a,b] \rightarrow \R: f\text{ is bounded}\}$$
	\end{defn}
	
	\begin{defn} \ld{00000} 
		Let $a,b \in \R$. 
		\begin{itemize}
			\item Let $n \in \N$ and $(x_j)_{j \in [n]_0} \in \R^{n+1}$. Then $(x_j)_{j \in [n]_0}$ is said to be a \tbf{Darboux partition of $[a,b]$} if $(x_j)_{j \in [n]_0}$ is strictly increasing, $x_0 = a$ and $x_n = b$. 
			\item We define $\Part_D([a,b]) \defeq \bigg\{\MP \in \bigcup\limits_{n \in \N} \R^{n+1} : \text{$\MP$ is a Darboux partition of $[a,b]$}\bigg\}$ \tcr{rewrite rest of section using this notation}
		\end{itemize}
	\end{defn}
	
	\begin{note}\
		\begin{itemize}
			\item We emphasize that the definition of a ``partition" of $[a,b]$ here is different than the set-theoretic definition which refers to a collection of disjoint subsets of $[a,b]$ which cover all of $[a,b]$. 
			\item We recall the definition of $M^f_{s,t}, m^f_{s,t}$ from \rex{} \tcr{ref prev def here}  
		\end{itemize}
	\end{note}
	
	\begin{defn} \ld{00000} 
		Let $f \in \Bnd([a,b], \R)$ and $\MP \in \Part_D([a,b])$. Write $\MP = (x_j)_{j \in [n]_0}$. 
		We define the 
		\begin{itemize}
			\item \textbf{upper Darboux sum} of $f$ with respect to $\MP$, denoted $U_\MP f$, by 
			$$U_\MP f \defeq \sum\limits_{i=1}^n M^f_{x_{j-1}, x_j} (x_j - x_{j-1})$$ 
			\item \textbf{lower Darboux sum} of $f$ with respect to $\MP$, denoted $L_\MP f$, by
			$$L_\MP f \defeq \sum\limits_{i=1}^n m^f_{x_{j-1}, x_j} (x_j - x_{j-1})$$ 
		\end{itemize}
	\end{defn}
	
	\begin{ex} \lex{00000} 
		Let $f \in \Bnd([a,b], \R)$ and $\MP$ a partition of $[a,b]$. Then 
		$$\bigg[\inf_{x \in [a,b]} f(x) \bigg] (b-a) \leq L_\MP f \leq U_\MP f \leq \bigg[ \sup_{x \in [a,b]} f(x)  \bigg] (b-a)$$
	\end{ex}
	
	\begin{proof}
		Clear.
	\end{proof}
	
	\begin{defn}
		Let $\MP \in \Part_D([a,b])$. Write $\MP = (x_j)_{j \in [m]_0}$.
		\begin{itemize}
			\item Let $\MP' \in \Part_D([a,b])$. Write $\MP' = (x'_k)_{k \in [n]_0}$. Then $\MP'$ is said to be a \tbf{refinement of $\MP$}, denoted $\MP < \MP'$, if $\{x_j\}_{j \in [m]_0} \subset \{x'_k\}_{k \in [n]_0}$.
			\item We define the \tbf{mesh of $\MP$}, denoted $\mesh(\MP) \in R_+$, by 
			$$\mesh(\MP) \defeq \max\limits_{j \in m_0}$$
		\end{itemize}
		
	\end{defn}
	
	\begin{ex} \lex{00000} 
		Let $f \in \Bnd([a,b], \R)$ and $\MP$, $\MP'$ partitions of $[a,b]$. If $\MP \leq \MP'$, then 
		\begin{enumerate}
			\item $U_{\MP'} f \leq U_{\MP} f$
			\item $L_{\MP} f \leq L_{\MP'} f$
		\end{enumerate}
	\end{ex}
	
	\begin{proof} \
		\begin{enumerate}
			\item Assume that $\MP = \{x_0, \cdots, x_n\}$ and $\MP' = \MP \cup \{x'\}$. Then there exists $j \in \{1, \cdots, n\}$ such that $x_{j-1} < x' < x_j$. Define $$M'_1 = \sup_{x \in [x_{j-1}, x']} f(x), \hspace{.4cm} M'_2 = \sup_{x \in [x', x_{j}]} f(x)$$
			Since $[x_{j-1}, x'], [x', x_j] \subset [x_{j-1}, x_j]$, we have that $M'_1, M'_2 \leq  M^f_j$. Then 
			\begin{align*}
				U_{P'}f 
				&= \sum\limits_{i =1}^{j-1} M^f_i(x_i - x_{i-1}) + M'_1(x' - x_{j-1}) + M'_2(x_j - x') + \sum\limits_{i =j+1}^n M^f_i(x_i - x_{i-1})  \\
				&\leq   \sum\limits_{i =1}^{n} M^f_i(x_i - x_{i-1}) \\
				&= U_P f 
			\end{align*}
			By induction, this is true for general partitions $P \leq \MP'$.
			\item Similar to $(1)$.
		\end{enumerate}
	\end{proof}
	
	\begin{ex} \lex{00000} 
		Let $f, g \in \Bnd([a,b], \R)$ and $\MP = \{x_0, \cdots, x_n\}$ a partition of $[a,b]$. Then 
		\begin{enumerate}
			\item $U_\MP(f+g) \leq U_\MP f + U_\MP g$
			\item $L_\MP(f+g) \geq L_\MP f + L_\MP g$
		\end{enumerate}
	\end{ex}
	
	\begin{proof} \
		\begin{enumerate}
			\item For each $i \in \{1, \cdots, n\}$, $M^{f+g}_i \leq M^f_i + M^g_i$. So 
			\begin{align*}
				U_\MP (f+g) 
				&= \sum\limits_{i=1}^n M^{f+g}_i (x_i - x_{i-1}) \\
				& \leq \sum\limits_{i=1}^n (M^f_i + M^g_i)(x_i - x_{i-1}) \\
				&= \sum\limits_{i=1}^n M^f_i(x_i - x_{i-1}) + \sum\limits_{i=1}^n M^g_i(x_i - x_{i-1}) \\
				&= U_\MP f + U_\MP g
			\end{align*}
			\item Similar to $(1)$
		\end{enumerate}
	\end{proof}
	
	\begin{ex} \lex{00000} 
		Let $f \in \Bnd([a,b], \R)$ and $\MP = \{x_0, \cdots, x_n\}$ a partition of $[a,b]$. Then 
		\begin{enumerate}
			\item $U_\MP (-f) = -L_P f$
			\item $L_\MP (-f) = - U_\MP f$
		\end{enumerate}
	\end{ex}
	
	\begin{proof}\
		\begin{enumerate}
			\item Since for $i \in \{1, \cdots, n\}$, $M^{-f}_i = -m^f_i$ we see that 
			\begin{align*}
				U_\MP(-f) 
				&= \sum\limits_{i=1}^n M^{-f}_i(x_i - x_{i-1}) \\
				&= - \sum\limits_{i=1}^n m^f_i(x_i - x_{i-1}) \\
				&= - L_\MP f
			\end{align*}
			\item Similar to $(1)$.
		\end{enumerate}
	\end{proof}
	
	\begin{ex} \lex{00000} 
		Let $f \in \Bnd([a,b], \R)$, $c >0$ and $\MP = \{x_0, \cdots, x_n\}$ a partition of $[a,b]$. Then 
		\begin{enumerate}
			\item $U_\MP (cf) = c U_\MP f$ 
			\item $L_\MP (cf) = c L_\MP f $
		\end{enumerate}
	\end{ex}
	
	\begin{proof}\
		\begin{enumerate}
			\item Since for $i \in \{1, \cdots, n\}$, $M^{cf}_i = cM^f_i$, we see that  
			\begin{align*}
				U_\MP (cf) 
				&= \sum\limits_{i=1}^n M^{cf}_i (x_i - x_{i-1}) \\
				&= c \sum\limits_{i=1}^n M^f_i (x_i - x_{i-1}) \\
				&= c U_\MP f
			\end{align*}
			\item Similar to $(1)$
		\end{enumerate}
	\end{proof}
	
	\begin{defn} \ld{00000} 
		Let $f \in \Bnd([a,b], \R)$. We define the \textbf{upper Darboux integral} of $f$, denoted $U f$, by
		$$Uf \defeq \inf \{U_\MP f: \MP \in \Part_D([a,b]) \}$$
		and we define the \textbf{lower Darboux integral} of $f$, denoted $L f$, by 
		$$Lf \defeq \sup \{L_\MP f: \MP \in \Part_D([a,b]) \}$$ 
	\end{defn}
	
	\begin{ex} \lex{00000} 
		Let $f \in \Bnd([a,b], \R)$. Then $$ \bigg[\inf_{x \in [a,b]} f(x) \bigg] (b-a) \leq Lf \leq Uf \leq \bigg[\sup_{x \in [a,b]} f(x) \bigg] (b-a)$$
	\end{ex}
	
	\begin{proof}
		Clearly 
		$$\bigg[\inf_{x \in [a,b]} f(x) \bigg] (b-a) \leq Lf \hspace{.2cm } \text{ and } \quad  Uf \leq \bigg[\sup_{x \in [a,b]} f(x) \bigg] (b-a)$$ 
		Let $\ep >0$. Then there exist partitions $\MP_1$ and $\MP_2$ of $[a,b]$ such that $U_{\MP_1} f < U f + \ep/2$ and $L_{\MP_2} f > Lf - \ep/2 $. Define $\MP = \MP_1 \cup \MP_2$. Then 
		\begin{align*}
			Uf 
			& \geq U_{\MP_1}  - \ep/2 \\
			&> U_\MP f - \ep/2 \\
			&\geq L_\MP f -\ep/2 \\
			&\geq L_{\MP_2} f -\ep/2 \\
			&> L f - \ep
		\end{align*} 
		Since $\ep >0$ is arbitrary, we have that $Uf \geq Lf$.
	\end{proof}
	
	\begin{ex} \lex{00000} 
		Let $f, g \in \Bnd([a,b], \R)$. Then 
		\begin{enumerate}
			\item $U(f+g) \leq U f + U g$
			\item $L(f+g) \geq L f + L g$
		\end{enumerate}
	\end{ex}
	
	\begin{proof}\
		\begin{enumerate}
			\item Let $\ep >0$. Then there exists a partitions $\MP_1$ of $[a,b]$ such that $U_{\MP_1} f < U f + \ep/2$ and $U_{\MP_2} g < U f + \ep/2$. Define $\MP = \MP_1 \cup \MP_2$. Then 
			\begin{align*}
				U_\MP(f +g) 
				&\leq U_\MP f + U_\MP g \\
				&\leq U_{\MP_1} f + U_{\MP_2} g \\
				&< U f + \ep/2 +  U g + \ep/2 \\
				&= Uf + Ug + \ep
			\end{align*}
			Since $\ep >0$ is arbitrary, $U_\MP(f +g)  \leq Uf + Ug$.
			\item Similar to $(1)$.
		\end{enumerate}
	\end{proof}
	
	\begin{ex} \lex{00000} 
		Let $f \in \Bnd([a,b], \R)$. Then 
		\begin{enumerate}
			\item $U(-f) = -Lf$
			\item $L(-f) = -Uf$
		\end{enumerate}
	\end{ex}
	
	\begin{proof} \
		\begin{enumerate}
			\item Using a previous exercise, we have that
			\begin{align*}
				U(-f)
				&= \inf \{U_\MP(-f): \MP \text{ is a partition of } [a,b]\} \\
				&= \inf \{ -L_\MP f: \MP \text{ is a partition of } [a,b]\} \\
				&= - \sup \{L_\MP f: \MP \text{ is a partition of } [a,b]\} \\
				&= - Lf
			\end{align*}
			\item Similar to $(1)$
		\end{enumerate}
	\end{proof}
	
	\begin{ex} \lex{00000} 
		Let $f \in \Bnd([a,b], \R)$ and $c \geq 0$. Then 
		\begin{enumerate}
			\item $U (cf) = c U f$ 
			\item $L (cf) = c L f $
		\end{enumerate}
	\end{ex}
	
	\begin{proof}\
		\begin{enumerate}
			\item Using a previous exercise, we have that
			\begin{align*}
				U(cf) 
				&= \inf \{U_\MP (cf): \MP \text{ is a partition of } [a,b]\} \\
				&= \inf \{c U_\MP f: \MP \text{ is a partition of } [a,b]\} \\
				&= c \inf \{U_\MP f: \MP \text{ is a partition of } [a,b]\} \\
				&= c Uf
			\end{align*}
			\item Similar to $(1)$
		\end{enumerate}
	\end{proof}
	
	\begin{defn} \ld{00000} \tbf{Darboux Integral:} \\ 
		\begin{itemize}
			\item Let $f \in \Bnd([a,b], \R)$.Then $f$ is said to be \textbf{Darboux integrable} if $Uf = Lf$. 
			\item We define $\Darb([a,b], \R) \subset \Bnd([a,b], \R)$ by 
			$$\Darb([a,b], \R) \defeq \{f \in \Bnd([a,b], \R): \text{$f$ is Darboux integrable}\}.$$
			\item Let $f \in \Darb([a,b])$. We define the \textbf{Darboux integral} of $f$, denoted
			$$I^D(f), \quad \int f, \quad  \int_a^b f, \quad \text{or} \quad \int_a^b f(x) dx,$$  
			by
			$$I^D_{[a,b]} (f) \defeq Uf.$$ 
		\end{itemize}
	\end{defn}
	
	\begin{note}
		When the context is clear, we may write $I(f)$ in place of $I^D(f)$.
	\end{note}
	
	\begin{ex} \lex{00000} 
		Let $f \in \Bnd([a,b], \R)$. Then $f \in \Darb([a,b], \R)$ iff for each $\ep >0$, there exists a partition $\MP$ of $[a,b]$ such that $U_\MP f - L_\MP f < \ep$.
	\end{ex}
	
	\begin{proof}
		Suppose that $f \in \Darb([a,b], \R)$. Let $\ep >0$. Then there exist partions $\MP_1$, $\MP_2$ of $[a,b]$ such that $U_{\MP_1} f < Uf + \ep/2$ and $L_{\MP_2} f > Lf - \ep/2$. Define $\MP = \MP_1 \cup \MP_2$. Then $U_\MP f \leq U_{\MP_1} f$ and $L_P f \geq L_{\MP_2}f$. So  
		\begin{align*}
			U_\MP f - L_\MP f 
			&< Uf - L f + \ep \\
			&= \ep
		\end{align*}  
		Conversely, suppose that for each $\ep >0$, there exists a partition $\MP$ of $[a,b]$ such that $U_\MP f - L_\MP f < \ep$. For the sake of contradiction, suppose that $Uf - Lf > 0$. Choose $\ep = Uf - Lf$. Then there exists a partition $\MP$ of $[a,b]$ such that $U_\MP f - L_\MP f < \ep$. Since $Uf \leq U_\MP f$ and $Lf \geq L_\MP f$, we have that 
		\begin{align*}
			\ep 
			&> U_\MP f - L_\MP f \\
			&\geq Uf - Lf \\
			&= \ep
		\end{align*} 
		which is a contradiction. Hence $Uf = Lf$ and $f \in \Darb([a,b], \R)$.
	\end{proof}
	
	\begin{ex} \lex{00000} 
		Let $f, g \in \Darb([a,b], \R)$. Then $f+g \in \Darb([a,b], \R)$ and $$\int (f+g) = \int f + \int g$$
	\end{ex}
	
	\begin{proof}
		Clearly $f+g \in \Bnd([a,b], \R)$. Using some previous results, we have that 
		\begin{align*}
			\int f + \int g 
			&= Lf + Lg \\
			&\leq L(f+g) \\
			&\leq U(f+g) \\
			&\leq Uf + Ug \\
			&= \int f + \int g
		\end{align*}
		So $U(f+g) = L(f+g) = \int f + \int g$. Therefore $f+g \in \Darb([a,b], \R)$ and $$\int (f+g) = \int f + \int g$$.
	\end{proof}
	
	\begin{ex} \lex{00000} 
		Let $f \in \Darb([a,b], \R)$ and $c \in \R$. Then $cf \in \Darb([a,b], \R)$ and $$\int (cf) = c \int f $$
	\end{ex}
	
	\begin{proof}
		Clearly $cf \in \Bnd([a,b], \R)$. If $c \geq 0$, then 
		\begin{align*}
			L (cf) 
			&= c Lf \\
			&= c \int f \\
			&= c Uf \\
			&= U(cf)
		\end{align*} 
		So $$L(cf) = U(cf) = c \int f$$ If $c <0$, then 
		\begin{align*}
			L (cf) 
			&= L(-|c| f) \\
			&= -U(|c|f) \\
			&= - |c|Uf \\
			&= c \int f \\
			&= -|c| L f \\
			&= - L(|c|f) \\
			&= U(-|c|f) \\
			&= U (c f) \\
		\end{align*} 
		So $$L(cf) = U(cf) = c \int f$$ Therefore $cf \in \Darb([a,b], \R)$ and $$\int (cf) = c \int f$$
	\end{proof}
	
	\begin{cor}
		We have that $\Darb([a,b], \R)$ is a vector space and the map $I: \Darb([a,b], \R) \rightarrow \R$ given by $If = \int f$ is linear.
	\end{cor}
	
	\begin{proof}
		Clear.
	\end{proof}
	
	\begin{ex} \lex{00000} 
		Let $f:[a,b] \rightarrow \R$. If $f$ is continuous, then $f \in \Darb([a,b], \R)$. 
	\end{ex}
	
	\begin{proof}
		Suppose that $f$ is continuous. Then $f$ is uniformly continuous. Let $\ep >0$. Uniform continuity implies that there exists $\del > 0$ such that for each $x, y \in [a,b]$, $|x-y| < \del$ implies that $|f(x) - f(Y)| < \ep/(b-a)$. Choose $n \in \N$ such that $(b-a)/n < \del$. For $i \in \{0, \cdots, n\}$, define $x_i \defeq a + i(b-a)/n$. Put $\MP = \{x_0, \cdots, x_n\}$. Continuity implies that for each $i \in \{1, \cdots, n\}$, there exists $x_i^M, x_i^m \in [x_{i-1}, x_i]$ such that $f(x_i^M) = M_i^f$ and $f(x_i^m) = m_i^f$. Then 
		\begin{align*}
			U_\MP f - L_\MP f 
			&= \sum\limits_{i=1}^n M^f_i(x_i - x_{i-1}) - \sum\limits_{i=1}^n m^f_i(x_i - x_{i-1}) \\
			&= \sum\limits_{i=1}^n (M^f_i - m^f_i)(x_i - x_{i-1}) \\
			&= \sum\limits_{i=1}^n [f(x_i^M) - f(x_i^m)](x_i - x_{i-1}) \\
			& < \sum\limits_{i=1}^n \frac{\ep}{b-a}(x_i - x_{i-1}) \\
			&= \ep 
		\end{align*}
		So for each $\ep >0$, there exists a partition $\MP$ of $[a,b]$ such that $U_\MP f - L_\MP f < \ep$. Hence $f \in \Darb([a,b], \R)$.
	\end{proof}
	
	\begin{ex} \lex{00000} 
		Let $f:[a,b] \rightarrow \R$. If $f$ is monotonic, then $f \in \Darb([a,b], \R)$.
	\end{ex}
	
	\begin{proof}
		Suppose that $f$ is increasing. Let $\ep >0$. Choose $n \in \N$ such that $(b-a)[f(b) - f(a)] /n < \ep $. For $i \in \{0, \cdots, n\}$, define $x_i \defeq a + i(b-a)/n$. Put $\MP = \{x_0, \cdots, x_n\}$. Then 
		\begin{align*}
			U_\MP f - L_\MP f 
			&= \sum\limits_{i=1}^n M^f_i(x_i - x_{i-1}) - \sum\limits_{i=1}^n m^f_i(x_i - x_{i-1}) \\
			&= \sum\limits_{i=1}^n (M^f_i - m^f_i)(x_i - x_{i-1}) \\
			&= \frac{b-a}{n} \sum\limits_{i=1}^n [f(x_i) - f(x_{i-1})]  \\
			&= \frac{b-a}{n} [f(b) - f(a)] \\
			& < \ep
		\end{align*}
		So for each $\ep >0$, there exists a partition $\MP$ of $[a,b]$ such that $U_\MP f - L_\MP f < \ep$. Hence $f \in \Darb([a,b], \R)$. The case is similar if $f$ is decreasing. 
	\end{proof}
	
	\begin{ex}
		Let $f \in \Darb([a,b], \R)$. Then for each $c \in [a,b]$, $f|_{[a,c]} \in \Darb([a,c], \R)$.
	\end{ex}
	
	\begin{proof}
		\tcr{make sure indices are correct}
		Let $c \in [a,b]$. Define $g:[a,c] \rightarrow \R$ by $g \defeq f|_{[a,c]}$. Let $\ep > 0$. Since $f \in \Darb([a,b], \R)$, \tcr{ref prev ex} implies that there exists a partition $\MP$ of $[a,b]$ such that $U_{\MP}f - L_{\MP}f < \ep$. Define $\MP_c$ by $\MP_c \defeq \MP \cup \{c\} \setminus (c,b]$. Then $\MP_c$ is a partition of $[a,c]$. Write $\MP = (x_j)_{j \in [n]_0}$ and $\MP = (y_j)_{j \in [l]_0}$. Then $l-1 = \max_{j \in [n]_0: x_j < c}$ and for each $j \in [l-1]_0$, $x_j = y_j$, 
		\begin{align*}
			M^g_j
			& = \sup_{x \in [y_{j-1}, y_j]} g(x) \\
			& = \sup_{x \in [x_{j-1}, x_j]} f|_{[a,c]}(x) \\
			& = \sup_{x \in [x_{j-1}, x_j]} f(x) \\
			& = M^f_j,
		\end{align*}
		and 
		\begin{align*}
			m^g_j
			& = \inf_{x \in [y_{j-1}, y_j]} g(x) \\
			& = \inf_{x \in [x_{j-1}, x_j]} f|_{[a,c]}(x) \\
			& = \inf_{x \in [x_{j-1}, x_j]} f(x) \\
			& = m^f_j.
		\end{align*}
		Since $y_l = c \in (x_{l-1}, x_l]$, we have that $y_l - y_{l-1} < x_l - x_{l-1}$,
		\begin{align*}
			M^g_l
			& = \sup_{x \in [y_{l-1}, y_l]} g(x) \\
			& = \sup_{x \in [x_{l-1}, c]} f|_{[a,c]}(x) \\
			& = \sup_{x \in [x_{l-1}, c]} f(x) \\
			& \leq \sup_{x \in [x_{l-1}, x_l]} f(x) \\
			& = M^f_l,
		\end{align*}
		and
		\begin{align*}
			m^g_l
			& = \inf_{x \in [y_{l-1}, y_l]} g(x) \\
			& = \inf_{x \in [x_{l-1}, c]} f|_{[a,c]}(x) \\
			& = \inf_{x \in [x_{l-1}, c]} f(x) \\
			& \geq \inf_{x \in [x_{l-1}, x_l]} f(x) \\
			& = m^f_l.
		\end{align*}
		Therefore
		\begin{align*}
			U_{\MP_c}g - L_{\MP_c} g
			& = \sum\limits_{j \in [l]}^l M^g_j (y_j - y_{j-1})  + \sum\limits_{j \in [l]} m^g_j (y_j - y_{j-1}) \\
			& = \sum\limits_{j \in [l]}  (M^g_j - m^g_j) (y_j - y_{j-1}) \\
			& = \sum\limits_{j \in [l-1]}  (M^g_j - m^g_j) (y_j - y_{j-1}) + (M^g_l - m^g_l) (y_l - y_{l-1}) \\
			& \leq \sum\limits_{j \in [l-1]}  (M^f_j - m^f_j) (x_j - x_{j-1}) + (M^g_l - m^g_l) (x_l - x_{l-1}) \\
			& = \sum\limits_{j \in [l]}  (M^f_j - m^f_j) (x_j - x_{j-1}) \\
			& \leq \sum\limits_{j \in [n]}  (M^f_j - m^f_j) (x_j - x_{j-1}) \\
			& = \sum\limits_{j \in [n]}  (M^f_j) (x_j - x_{j-1})  - \sum\limits_{j \in [n]}  (m^f_j) (x_j - x_{j-1}) \\
			& = U_{\MP} f - L_{\MP} f \\
			& < \ep. 
		\end{align*}  
		Since $\ep > 0$ is arbitrary, we have that for each $\ep > 0$, there exists a partition $\MP_c$ of $[a,c]$ such that such that $U_{\MP_c} g - L_{\MP_c} g < \ep$. \tcr{A prev ex} implies that 
		\begin{align*}
			f|_{[a,c]}
			& = g \\
			& \in \Darb([a,b], \R).
		\end{align*} 
		Since $c \in [a,b]$ is arbitrary, we have that for each $c \in [a,b]$, $f|_{[a,c]} \in \Darb([a,b], \R)$.
 	\end{proof}
 	
	\begin{ex} \lex{00000} 
		Define $\chi_{\Q}:[0,1] \rightarrow \R$ by $$\chi_{\Q}(x) = \begin{cases}
			1 & x \in \Q \\
			0 & x \not \in \Q
		\end{cases}$$
		Then $\chi_\Q \not \in \Darb([a,b], \R)$.
	\end{ex}
	
	\begin{proof}
		Let $\MP = \{x_0, \cdots, x_n\}$ be a partition of $[0,1]$. Then for each $i \in \{1, \cdots, n\}$, $M^{\chi_{\Q}}_i = 1$ and $m^{\chi_{\Q}}_i = 0$. So $U_\MP \chi_\Q = 1$ and $L_\MP \chi_\Q = 0$. Since $\MP$ is arbitrary, we have that $U \chi_\Q = 1$ and $L \chi_\Q = 0$
	\end{proof}
	
	
	
	
	
	
	
	
	
	
	
	
	
	
	
	
	
	
	
	
	
	
	
	
	
	
	
	
	
	
	
	
	
	
	
	
	
	
	
	
	
	
	
	
	
	
	
	
	\newpage
	\section{The Riemann Integral}
	
	\begin{defn} \ld{00000} 
		Let $a,b \in \R$. 
		\begin{itemize}
			\item Let $\MP \in \Part_D([a,b])$. Write $\MP = (x_j)_{j \in [n]_0}$.
			\begin{itemize}
				\item Let $\MT \in \R^n$. Write $\MT = (t_k)_{k \in [n]}$. Then
				\begin{itemize}
					\item $\MT$ is said to be a \tbf{tag of $\MP$} if for each $k \in [n]$, $t_k \in [x_k, x_{k-1}]$,
					\item $(\MP, \MT)$ is said to be a \tbf{tagged partition of $[a,b]$}.
				\end{itemize}
				\item We define $\Tag_{\MP} \defeq \{\MT \in \R^n: \text{$\MT$ is a tag of $\MP$}\}$.
			\end{itemize}
			\item We define $\Part_R([a,b]) \defeq \coprod\limits_{\MP \in \Part_D([a,b])} \Tag_{\MP}$
		\end{itemize}
	\end{defn}
	
	\begin{defn} \ld{00000} 
		Let $f \in \Bnd([a,b], \R)$ and $(\MP, \MT) \in \Part_R([a,b])$. Write $\MP = (x_j)_{j \in [n]_0}$ and $\MT = (t_k)_{k \in [n]})$. We define the \textbf{Riemann sum} of $f$ with respect to $(\MP, \MT)$, denoted $S^f(\MP, \MT) f$, by 
		$$S^f(\MP, \MT) \defeq \sum\limits_{j=1}^n f(t_j) (x_j - x_{j-1}).$$ 
	\end{defn}
	
	\begin{ex} \lex{00000} 
		Let $f \in \Bnd([a,b], \R)$ and $(\MP, \MT) \in \Part([a,b])$. Write $\MP = ((x_j)_{j \in [n]_0}$ and $(t_k)_{k \in [n]})$. Define $\MP_D \in \Part_D([a,b])$ by $\MP_d \defeq (x_j)_{j \in [n]_0}$. Then 
		$$L^f(\MP) \leq S^f(\MP, \MT) \leq U^f(\MP)$$ 
	\end{ex}
	
	\begin{proof}
		Clear.
	\end{proof}
	
	\begin{ex} \lex{00000} 
		Let $f, g \in \Bnd([a,b], \R)$ and $(\MP, \MT) \in \Part([a,b])$. Then $S^{f+g}(\MP, \MT) = S^f(\MP, \MT) + S^g(\MP, \MT)$.
	\end{ex}
	
	\begin{proof} 
		Write $\MP = (x_j)_{j \in [n]_0}$ and $\MT = (t_k)_{k \in [n]})$. Then 
		\begin{align*}
			S^{f+g}(\MP, \MT)
			& = \sum\limits_{j=1}^n [f+g](t_j) (x_j - x_{j-1}) \\
			& = \sum\limits_{j=1}^n [f(t_j) + g(x_j)] (x_j - x_{j-1}) \\
			& = \sum\limits_{j=1}^n [f(t_j) (x_j - x_{j-1}) + g(x_j) (x_j - x_{j-1})] \\
			& = \sum\limits_{j=1}^n f(t_j) (x_j - x_{j-1}) +  \sum\limits_{j=1}^n g(x_j) (x_j - x_{j-1}) \\
			& = S^f(\MP, \MT) + S^g(\MP, \MT).
		\end{align*}
	\end{proof}
	
	\begin{ex} \lex{00000} 
		Let $f \in \Bnd([a,b], \R)$, $\lam \in \R$ and $(\MP, \MT) \in \Part([a,b])$. Then $S^{\lam f}(\MP, \MT) = \lam S^f(\MP, \MT)$.
	\end{ex}
	
	\begin{proof}
		Write $\MP = ((x_j)_{j \in [n]_0}$ and $(t_k)_{k \in [n]})$. Then 
		\begin{align*}
			S^{\lam f}(\MP, \MT) 
			& = \sum\limits_{j=1}^n [\lam f](t_j) (x_j - x_{j-1}) \\
			& = \sum\limits_{j=1}^n \lam f(t_j) (x_j - x_{j-1}) \\
			& = \lam \sum\limits_{j=1}^n f(t_j) (x_j - x_{j-1}) \\
			& = \lam S^f(\MP, \MT).
		\end{align*}
	\end{proof}
	
	\begin{defn} \ld{00000} 
		Let $a,b \in \R$ and $(\MP, \MT), (\MP', \MT') \in \Part_R([a,b])$. Write $\MT=\{s_j\}_{j \in [m]}$ and $\MT' = \{t_k\}_{k \in [n]}$. Then $(\MP', \MT')$ is said to be a \tbf{refinement of $(\MP, \MT)$}, denoted $(\MP, \MT) \leq (\MP', \MT')$, if  
		\begin{enumerate}
			\item $\MP \leq \MP'$,
			\item $\{s_j\}_{j \in [m]} \subset \{t_k\}_{k \in [n]}$.
		\end{enumerate} 
	\end{defn}
	
	\begin{defn}
		\begin{itemize}
			Let $f:[a,b] \rightarrow \R$. 
			\begin{itemize}
				\item Let $I \in \R$. Then $I$ is said to be \tbf{a Riemann integral of $f$} if for each $\ep > 0$, there exists $\del > 0$ such that for each $(\MP,\MT) \in \Part_R([a,b])$, $\mesh(\MP) < \del$ implies that $| S^f(\MP, \MT) - I| < \ep$.  
				\item We define $\IntSet(f) \subset \R$ by $\IntSet(f) \defeq \{I \in \R: \text{$I$ is a Riemann integral of $f$}\}$.
				\item Then $f$ is said to be \tbf{Riemann integrable} if $\IntSet(f) \neq \varnothing$.
			\end{itemize}
			\item We define $\Riem([a,b]) \defeq \{f \in \R^{[a,b]}: \text{$f$ is Riemann integrable}\}$.
		\end{itemize}
	\end{defn}
	
	\begin{ex}
		We have that $\Riem([a,b]) \subset \Bnd([a,b], \R)$.
	\end{ex}
	
	\begin{proof}
		\tcr{FINISH!!!!}
	\end{proof}
	
	\begin{ex}
		Let $f: [a,b] \rightarrow \R$. Then $f \in \Riem([a,b])$ iff there exists $I \in \R$ such that $\IntSet(f) = \{I\}$.
	\end{ex}
	
	\begin{proof}\
		\begin{itemize}
			\item $(\implies)$: \\
			\begin{itemize}
				\item If $a=b$, \tcr{FINISH!!!}
				\item Suppose that $a \neq b$. Suppose that $f \in \Riem([a,b])$. Then $\IntSet(f) \neq \varnothing$ and there exists $I \in \R$ such that $I \in \IntSet(f)$. Let $I' \in \IntSet(f)$. Let $\ep > 0$. Then $\ep/2 > 0$. Since $I \in \IntSet(f)$, there exists $\del_1 > 0$ such that for each $(\MP,\MT) \in \Part_R([a,b])$, $\mesh(\MP) < \del_1$ implies that $| S^f(\MP, \MT) - I| < \ep/2$. Since $I' \in \IntSet(f)$, there exists $\del_2 > 0$ such that for each $(\MP,\MT) \in \Part_R([a,b])$, $\mesh(\MP) < \del_2$ implies that $| S^f(\MP, \MT) - I'| < \ep/2$. Define $\del \defeq \min \{\del_1, \del_2\}$. There exists $n \in \N$ such that $(b-a)/n < \del$. For each $j \in [n]_0$, define $x_j \in [a,b]$ by $x_j \defeq a + j(b-a)/n$. Define $(\MP, \MT) \in \Part_R([a,b])$ by $\MP \defeq (x_j)_{j \in [n]_0}$ and $\MT \defeq (x_j)_{j \in [n]}$. Since for each $j \in [n]$, 
				\begin{align*}
					x_j - x_{j-1} 
					& = \bigg[a + \frac{j(b-a)}{n} \bigg] - \bigg[a + \frac{(j-1)(b-a)}{n} \bigg] \\
					& = \frac{j(b-a)}{n} - \frac{(j-1)(b-a)}{n} \\
					& = [j - (j-1)]\frac{(b-a)}{n} \\
					& = \frac{(b-a)}{n} \\
					& < \del,
				\end{align*}
				we have that
				\begin{align*}
					\mesh(\MP)
					& = \max\limits_{j \in [n]} x_j - x_{j-1} \\
					& < \del. 
				\end{align*}
				Since $\del \leq \del_1$ and $\del \leq \del_2$, we have that for each $(\MP,\MT) \in \Part_R([a,b])$, $\mesh(\MP) < \del$ implies that 
				\begin{align*}
					|I - I'| 
					& \leq |I - S^f(\MP, \MT)| + | S^f(\MP, \MT) - I'| \\
					& < \frac{\ep}{2} + \frac{\ep}{2} \\
					& = \ep. 
				\end{align*} 
				Since $\ep > 0$ is arbitrary, we have that for each $\ep > 0$, $|I - I'| < \ep$. Thus $|I - I'| = 0$ and $I = I'$. Since $I' \in \IntSet(f)$ is arbitrary, we have that for each $I' \in \IntSet(f)$, $I' = I$. Hence
				\begin{align*}
					\{I\}
					& \subset \IntSet(f) \\
					& \subset \{I\}.
				\end{align*}
				Therefore $\IntSet(f) = \{I\}$. 
			\end{itemize}
			\item $(\impliedby)$: \\
			Suppose that there exists $I \in \R$ such that $\IntSet(f) = \{I\}$. Then $\IntSet(f) \neq \varnothing$ and $f \in \Riem([a,b])$. 
		\end{itemize}
	\end{proof}
	
	\begin{defn}
		Let $f \in \Riem([a,b])$. We define the \textbf{Riemann integral} of $f$, denoted
		$$I^R(f), \quad \int f, \quad  \int_a^b f, \quad \text{or} \quad \int_a^b f(x) dx,$$  
		to be the unique $I \in \R$ such that $\Riem(f) = \{I\}$.
	\end{defn}
	
	\begin{ex}
		
	\end{ex}
	
	\begin{ex}
		Let $a,b \in \R$. Suppose that $a \leq b$. Then 
		\begin{enumerate}
			\item $\Riem([a,b]) = \Darb([a,b])$
			\item for each $f \in \Riem([a,b])$, $I^R(f) = I^D(f)$. 
		\end{enumerate}
	\end{ex}
	
	\begin{proof}\
		\begin{enumerate}
			\item 
			\begin{itemize}
				\item $a =b$
				\item $a < b$  \tcr{FINISH!!!}
				\item Let $f \in \Riem([a,b])$ and $\ep > 0$. Then there exists $\del > 0$ such that for each $(\MP, \MT) \in \Part_R([a,b])$, $\mesh(\MP) < \del$ implies that $|I^R(f) - S^f(\MP, \MT)| < \ep/2$. 
				
				Define $\MP \in \Part_D([a,b])$ by $\MP \defeq (x_j)_{j \in [m]_0}$. Then 
				\begin{align*}
					\mesh(\MP)
					& = \max\limits_{j \in [m]} x_j - x_{j-1} \\
					& < \del. 
				\end{align*}
				Therefore
				\begin{align*}
					\sum\limits_{j \in [m]} \bigg( f(t_j) -  f(s_j)\bigg)  (x_j - x_{j-1})
					& = \sum\limits_{j \in [m]} \bigg[ f(t_j) (x_j - x_{j-1}) - f(s_j) (x_j - x_{j-1}) \bigg] \\
					& = \sum\limits_{j \in [m]} f(t_j) (x_j - x_{j-1})  - \sum\limits_{j \in [m]} f(s_j) (x_j - x_{j-1}) \\
					& = |I^R(f) - S^f(\MP, \MT)| \\
					& < \frac{\ep}{2}.
				\end{align*}
				
				For each $j \in [m]$, there exists $t_j,s_j \in [x_{j-1}, x_j]$ such that $M^f_{x_{j-1}, x_j} -\ep < f(t_j)$ and $f(s_j) < m^f_{x_{j-1}, x_j} + \ep$. Then there exist $n \in \N$ and $(y_k)_{k \in [n]_0}$ such that $\{x_j\}_{j \in [m]_0} \cup \{t_j\}_{j \in [m]} \cup \{s_j\}_{j \in [m]} = \{y_k\}_{k \in [n]_0}$. Define $\MP' \in \Part_D([a,b])$ by $\MP' \defeq \{y_k\}_{k \in [n]_0}$. Then $\MP \leq \MP'$. Define $\phi:[n] \rightarrow [m]$ by $\phi(k) \defeq \min \{j \in [m]: y_k \leq x_j\}$. \tcr{(maybe give some more details abou how $y_{\phi(k)} - y_{\phi(k) - 1} \leq x_j - x_{j-1}$)}. Thus for each $k \in [n]$, we have that $[y_{k-1}, y_k] \subset [x_{\phi(k) -1}, x_{\phi(k)}]$,
				\begin{align*}
					M^f_{y_{k-1}, y_k}
					& = \sup\limits_{x \in [y_{k-1}, y_k]} f(x) \\
					& \leq \sup\limits_{x \in [x_{\phi(k) -1}, x_{\phi(k)}]} f(x) \\
					& = M^f_{x_{\phi(k) -1}, x_{\phi(k)}},
				\end{align*}
				and 
				\begin{align*}
					m^f_{y_{k-1}, y_k}
					& = \inf\limits_{x \in [y_{k-1}, y_k]} f(x) \\
					& \geq \inf\limits_{x \in [x_{\phi(k) -1}, x_{\phi(k)}]} f(x) \\
					& = m^f_{x_{\phi(k) -1}, x_{\phi(k)}}.
				\end{align*}
				Furthermore, for each $j \in [m]$, $[x_{j-1}, x_j] = \bigcup\limits_{k \in \phi^{-1}j} [y_{k - 1}, y_k]$ and therefore $x_j - x_{j-1} = \sum\limits_{k \in \phi^{-1}(j)} (y_k - y_{k-1})$. 
				\begin{align*}
					U^f(\MP') - L^f(\MP')
					& = \bigg[ \sum\limits_{k \in [n]} M^f_{y_{k-1}, y_k} (y_k - y_{k-1}) \bigg] - \bigg[ \sum\limits_{k \in [n]} m^f_{y_{k-1}, y_k} (y_k - y_{k-1}) \bigg] \\
					& = \sum\limits_{k \in [n]} (M^f_{y_{k-1}, y_k} - m^f_{y_{k-1}, y_k}) (y_k - y_{k-1}) \\
					& = \sum\limits_{j \in [m]} \bigg[ \sum\limits_{k \in \phi^{-1}(j)} (M^f_{y_{k-1}, y_k} - m^f_{y_{k-1}, y_k}) (y_k - y_{k-1}) \bigg] \\
					& \leq \sum\limits_{j \in [m]} \bigg[ \sum\limits_{k \in \phi^{-1}(j)} (M^f_{x_{\phi(k) -1}, x_{\phi(k)}} - m^f_{x_{\phi(k) -1}, x_{\phi(k)}}) (y_k - y_{k-1}) \bigg] \\
					& = \sum\limits_{j \in [m]} \bigg[ \sum\limits_{k \in \phi^{-1}(j)} (M^f_{x_{j -1}, x_{j}} - m^f_{x_{j -1}, x_{j}}) (y_k - y_{k-1}) \bigg] \\
					& = \sum\limits_{j \in [m]} (M^f_{x_{j -1}, x_{j}} - m^f_{x_{j -1}, x_{j}})  \sum\limits_{k \in \phi^{-1}(j)} (y_k - y_{k-1})  \\
					& = \sum\limits_{j \in [m]} (M^f_{x_{j -1}, x_{j}} - m^f_{x_{j -1}, x_{j}})  (x_j - x_{j-1})  \\
					& < \sum\limits_{j \in [m]} \bigg( \bigg[f(t_j) + \frac{\ep}{4(b-a)} \bigg] - \bigg[f(s_j) - \frac{\ep}{4(b-a)} \bigg] \bigg)  (x_j - x_{j-1})  \\
					& = \sum\limits_{j \in [m]} \bigg( f(t_j) -  f(s_j) + \frac{\ep}{2(b-a)} \bigg)  (x_j - x_{j-1})  \\
					& = \sum\limits_{j \in [m]} \bigg[ \bigg( f(t_j) -  f(s_j)\bigg)  (x_j - x_{j-1}) +  \frac{\ep}{2(b-a)} (x_j - x_{j-1}) \bigg] \\
					& = \sum\limits_{j \in [m]} \bigg( f(t_j) -  f(s_j)\bigg)  (x_j - x_{j-1}) + \sum\limits_{j \in [m]}  \frac{\ep}{2(b-a)} (x_j - x_{j-1}) \\
					& = \sum\limits_{j \in [m]} \bigg( f(t_j) -  f(s_j)\bigg)  (x_j - x_{j-1}) + \frac{\ep}{2(b-a)} \sum\limits_{j \in [m]} (x_j - x_{j-1}) \\
					& = \sum\limits_{j \in [m]} \bigg( f(t_j) -  f(s_j)\bigg)  (x_j - x_{j-1}) + \frac{\ep (b-a)}{2(b-a)}  \\
					& = \sum\limits_{j \in [m]} \bigg( f(t_j) -  f(s_j)\bigg)  (x_j - x_{j-1}) + \frac{\ep}{2}  \\
					& < \frac{\ep}{2} + \frac{\ep}{2} \\
					& = \ep. 
				\end{align*}
				Since $\ep > 0$ is arbitrary, we have that for each $\ep > 0$, there exists $\MP' \in \Part_D([a,b])$ such that $U_{\MP'} - L_{\MP'} < \ep$. Thus $f \in \Darb([a,b])$. 
				\item Let $f \in \Darb([a,b])$. \tcr{FINISH!!!}
			\end{itemize}
			\item 
		\end{enumerate}
	\end{proof}
	
	
	
	
	
	
	
	
	
	
	
	
	
	
	
	
	
	
	
	
	
	
	
	
	
	
	
	
	
	
	
	
	
	
	
	
	
	
	
	
	
	
	\newpage
	\section{The Darboux-Stieltjes Integral}
	
	
	\begin{defn}
		Let $f \in \Bnd([a,b], \R)$, $G:[a,b] \rightarrow \R$ and $\MP \in \Part_D([a,b])$. Write $\MP = (x_j)_{j \in [n]_0}$. We define the 
		\begin{itemize}
			\item \textbf{upper Darboux-Stieltjes sum} of $f$ with respect to $\MP$, denoted $U_\MP (f,G)$, by 
			$$U_\MP f \defeq \sum\limits_{j \in [n]} M^f_j (G(x_j) - G(x_{j-1}))$$ 
			\item define the \textbf{lower Darboux-Stieltjes sum} of $f$ with respect to $\MP$, denoted $L_\MP (f,G)$, by
			$$L_\MP f \defeq \sum\limits_{j \in [n]} m^f_j (G(x_j) - G(x_{j-1}))$$ 
		\end{itemize}
	\end{defn}
	
	\begin{ex}
		Let $f \in \Bnd([a,b], \R)$, $G:[a,b] \rightarrow \R$. If $G$ is differentiable, then 
	\end{ex}
	
	
	
	\tcr{Oriented Darboux integral}
	
	\tcr{motivate path orientation with particle trajectory and reversing time}
	
	\begin{ex}
		Define $S:[a,b] \rightarrow [b,a]$ by $s(t) \defeq ta + (1-t) b$. Then 
		\begin{enumerate}
			\item $U_{\MP}(f,S) = - U_{\MP}$
			\item $L_{\MP}(f,S) = - L_{\MP}$
		\end{enumerate}
	\end{ex}
	
	\begin{proof}
		
	\end{proof}
	
	
	
	
	
	
	
	
	
	
	
	
	
	
	
	
	
	
	
	
	
	
	
	
	
	
	
	
	
	
	
	
	
	
	
	
	
	\section{Fundamental Theorem of Calculus}
	
	\tcr{try to motivate oriented traversal}
	Suppose we are given $x_0 \in [a,b]$. Then in many instances, it is natural to consider the cumulative area function of $f$ after $x_0$, which is denoted $F:[x_0 b] \rightarrow \R$, and defined by $$F(x) \defeq \int_{x_0}^x f.$$  
	For example this function can describe  
	\begin{itemize}
		\item various characteristics of physical systems like the total length of the trajectory a particle traverses between times $x_0$ and $x$, 
		\item various probabilistic outcomes like the chance of observing or measuring a value of an absolutely conitnuous random variable between $x_0$ and $x$.
	\end{itemize}
	Taking inspiriation from the particle example, 
	
	suppose that we start a timer at time $t =0$ and stop the timer at time $t=1$ and between these two times, the trajectory taken by a particle is described by the function $x:[0,1] \rightarrow \R$ (or perhaps it is easier to imagine in higher dimensions with $x:[0,1] \rightarrow \R^k$ for $k=2,3$). In other words, the position of the particle at time $t \in [0,1]$ as recorded by our timer is precisely $x(t)$. Now suppose the flow of time suddenly reverses once the particle reaches its terminal position $x(1)$ and we restart our timer. That is, at \tcr{FINISH!!!}
	
	\begin{align*}
		\int_a^b f \circ r 
	\end{align*}
	
	
	
	Then we wold observe the particle moving along the same unoriented trajectory as before, but in reverse. 
	We recall that \tcr{a prev ex} implies that for each $x \in [x_0, b]$,
	\begin{align*}
		\int_a^b f = \int_a^{x_0} f + \int_{x_0}^b f
	\end{align*}
	
	\begin{defn}
		Let $f \in \Darb([a,b], \R)$ and $x_0 \in [a,b]$. We define the \tbf{indefinite integral of $f$ from $x_0$}, denoted $I^f_{x_0}: [a,b] \rightarrow \R$, by 
		\[
		I^f_{x_0}(x) \defeq
		\begin{cases}
			I_{[x_0, x]} f, & x_0 \leq x \\
			-I_{[x, x_0]} f, & x_0 > x.
		\end{cases}
		\] 
	\end{defn}
	
	\begin{ex} \tbf{Fundamental Theorem of Calculus:} \\
		\begin{enumerate}
			\item Let $f \in \Darb([a,b], \R)$. Define $F: [a,b] \rightarrow \R$ by 
			$$F(x) \defeq \int_a^x f.$$ 
			Then for each $x_0 \in [a,b]$, $f$ is continuous at $x_0$ implies that $F$ is differentiable at $x_0$ and $F'(x_0) = f(x_0)$,  
			\item Let $F: [a,b] \rightarrow \R$. Suppose that $F$ is differentiable. Define $f: [a,b] \rightarrow \R$ by $f \defeq F'$. If $f \in \Darb([a,b], \R)$, then
			$$\int_a^b f = F(b) - F(a).$$
		\end{enumerate}
	\end{ex}
	
	\begin{proof}\
		\begin{enumerate}
			\item Let $x_0 \in (a,b)$. Since
			\begin{align*}
				\frac{F(x_0 + h) - F(x_0)}{h} 
				& = \frac{1}{h} \bigg[ \int_a^{x_0 + h} f  - \int_a^{x_0} f \bigg] \\
				& = \frac{1}{h} \int_{x_0}^{x_0 + h} f \quad \text{(by \tcr{make/ref ex})}
			\end{align*}
			and 
			\begin{align*}
				m^f_{x_0, x_0 + h} h
				& \leq  \int_{x_0}^{x_0 + h} f \\
				& \leq M^f_{x_0, x_0 + h} h,
			\end{align*}
			we have that 
			\begin{align*}
				m^f_{x_0, x_0 + h} 
				& = \frac{h}{h} m^f_{x_0, x_0 + h}  \\ 
				& \leq  \frac{F(x_0 + h) - F(x_0)}{h}  \\
				& \leq \frac{h}{h} M^f_{x_0, x_0 + h}  \\ 
				& = M^f_{x_0, x_0 + h}.
			\end{align*}
			Since $f$ is continuous at $x_0$, \rex{ex:analysis:real_nums:limits_funct:0010} implies that $\limh{0^+} M^f_{x_0, x_0 + h} = f(x_0)$ and $\limh{0^+} m^f_{x_0, x_0 + h} = f(x_0)$. \rex{} \tcr{make and ref monotonicity of limits (i.e. sandwhich) ex} implies that $\limh{0^+} h^{-1}[F(x_0 + h) - F(x_0)] = f(x_0)$. Since
			\begin{align*}
				F(x_0 - h) - F(x_0)
				& = - [F(x_0) - F(x_0 - h)] \\
				& = \bigg[ \int_a^{x_0 + h} f  - \int_a^{x_0} f \bigg] \\
				& = \int_{x_0}^{x_0 + h} f
			\end{align*}
			Similarly, \tcr{(add details)} $\limh{0^-} h^{-1}[F(x_0 + h) - F(x_0)] = f(x_0)$. Thus $\limh{0} h^{-1}[F(x_0 + h) - F(x_0)] = f(x_0)$, $F$ is differentiable at $x_0$ and $F'(x_0) = f(x_0)$.
			\item Suppose that $f \in \Darb([a,b], \R)$. Let $\ep > 0$. 
			\begin{itemize}
				\item Since  
				\begin{align*}
					\int_a^b f + \ep 
					& > \int_a^b f \\
					& = U f \\
					& = \inf \{U_{\MP} f: \text{$\MP$ is a partition of $[a,b]$}\},
				\end{align*}
				there exists a partion $\MP$ of $[a,b]$ such that $U_{\MP}f < 	\int_a^b f + \ep$. Write $\MP = (x_j)_{j \in [n]_0}$. Since $F$ is differentiable, the mean value theorem (\rex{}, \tcr{ref here}) implies that for each $j \in [n]$, there exists $c_j \in [x_{j-1}, x_j]$ such that 
				\begin{align*}
					f(c_j)
					& = F'(c_j) \\
					& = \frac{F(x_j) - F(x_{j-1})}{x_j - x_{j-1}}.
				\end{align*} 
				Therefore 
				\begin{align*}
					F(b) - F(a)
					& = F(x_n) - F(x_0) \\
					& = \sum\limits_{j \in [n]} F(x_j) - F(x_{j-1}) \\
					& = \sum\limits_{j \in [n]} f(c_j) (x_j - x_{j-1}) \\
					& \leq \sum\limits_{j \in [n]} M^f_{x_{j-1}, x_j} (x_j - x_{j-1}) \\
					& = U_{\MP} f \\
					& < \int_a^b f + \ep.
				\end{align*}
				\item Since  
				\begin{align*}
					\int_a^b f - \ep 
					& < \int_a^b f \\
					& = L f \\
					& = \sup \{L_{\MP} f: \text{$\MP$ is a partition of $[a,b]$}\},
				\end{align*}
				there exists a partion $\MP$ of $[a,b]$ such that $\int_a^b f - \ep < L_{\MP}f $. Write $\MP = (y_j)_{j \in [n]_0}$. Since $F$ is differentiable, the mean value theorem (\rex{}, \tcr{ref here}) implies that for each $j \in [n]$, there exists $d_j \in [y_{j-1}, y_j]$ such that 
				\begin{align*}
					f(d_j)
					& = F'(d_j) \\
					& = \frac{F(y_j) - F(y_{j-1})}{y_j - y_{j-1}}.
				\end{align*} 
				Therefore 
				\begin{align*}
					F(b) - F(a)
					& = F(y_n) - F(y_0) \\
					& = \sum\limits_{j \in [n]} F(y_j) - F(y_{j-1}) \\
					& = \sum\limits_{j \in [n]} f(d_j) (y_j - y_{j-1}) \\
					& \geq \sum\limits_{j \in [n]} m^f_{y_{j-1}, y_j} (y_j - y_{j-1}) \\
					& = L_{\MP} f \\
					& > \int_a^b f - \ep.
				\end{align*}
			\end{itemize}
			Since $\ep > 0$ is arbitrary, we have that 
			\begin{align*}
				\int_a^b f 
				& \leq F(b) - F(a) \\
				& \leq \int_a^b f.
			\end{align*} 
			Hence
			$$\int_a^b f = F(b) - F(a).$$ 
		\end{enumerate}
	\end{proof}
	
	
	
	
	
	
	
	
	
	
	
	
	
	
	
	
	
	
	
	
	
	
	
	
	
	
	
	
	
	
	
	
	
	
	
	
	
	
	
	
	
	
	
	\subsection{Logarithmic and Exponential Functions}
	
	\subsection{Logarithmic Function}
	
	\begin{defn}
		We define the \tbf{logarithmic function}, denoted $\log:  \rightarrow \R$, by 
		$$\log(x) \defeq \int_1^t \frac{1}{t} \dt. $$
	\end{defn}
	
	\begin{ex}
		
	\end{ex}
	
	\subsection{Exponential Function}
	
	\begin{note}
		\tcr{need to talk about summation, $l^1$, and ref here}
	\end{note}
	
	\begin{ex}
		For each $x \in \R$, $(x^n/n!)_{n \in \N_0} \in l^1(\N_0)$
	\end{ex}
	
	\begin{defn}
		We define the \tbf{exponential function}, denoted $\exp: \R \rightarrow \R$, by 
		$\exp(x) \defeq \sum\limits_{n \in \N_0}$
	\end{defn}
	
	\begin{note}
		When the context is clear, we may write $e^x$ in place of $\exp(x)$. 
	\end{note}
	
	\begin{ex}
		For each $x \in \R$, 
		$$e^x = \limn \bigg(1 + \frac{x}{n} \bigg)^n.$$
	\end{ex}
	
	\begin{proof} \tcr{need to define $\log$}.
		Let $x \in \R$. Then for each $n \in \N$, $n \geq |x|$ implies that
		\begin{align*}
			\log \bigg[ \bigg(1 + \frac{x}{n} \bigg)^n \bigg] 
			& = n \log \bigg(1 + \frac{x}{n} \bigg) \\
			& = \frac{\log (1 + xn^{-1} )}{n^{-1}}
		\end{align*}
		
		\tcr{FINISH!!!, Give more details}
		Define $f,g: (|x|, \infty) \rightarrow \R$ by $f(t) \defeq \log (1 + xt^{-1} )$ and $g(t) \defeq t^{-1}$. Then 
		\begin{align*}
			\frac{f'(t) }{g'(t)}
			& = \limt{\infty} \frac{-xt^{-2}(1 + xt^{-1})^{-1}}{-t^{-2}} \\
			& = \limt{\infty} x(1 + xt^{-1})^{-1} \\
			& = x, 
		\end{align*}
		$\limt f'(t) = 0$ and $\limt g'(t) = 0$. L'H\^opital's rule \tcr{(ref ex)} implies that
		\begin{align*}
			\limt{\infty} \frac{f(t)}{g(t)}
			& = \limt{\infty} \frac{f'(t)}{g'(t)} \\
			& = x.
		\end{align*}
		\rex{} \tcr{(ref ex)} then implies that 
		\begin{align*}
			\limn \log \bigg[ \bigg(1 + \frac{x}{n} \bigg)^n \bigg] 
			& = \limn \frac{\log (1 + x/n )}{1/n} \\
			& = \limn \frac{f(n)}{g(n)} \\
			& = \limn \frac{f'(n)}{g'(n)} \\
			& = x.
		\end{align*} 
		Continuity of $\exp$ then implies that 
		\begin{align*}
			\limn \bigg[ \bigg(1 + \frac{x}{n} \bigg)^n \bigg] 
			& = \limn \exp \circ \log \bigg[ \bigg(1 + \frac{x}{n} \bigg)^n \bigg] \\
			& = \limn \exp \bigg( \log \bigg[ \bigg(1 + \frac{x}{n} \bigg)^n \bigg] \bigg) \\
			& = \exp \bigg( \limn \log \bigg[ \bigg(1 + \frac{x}{n} \bigg)^n \bigg] \bigg) \\
			& = \exp(x).
		\end{align*}
		Since $x \in \R$ is arbitrary, we have that for each $x \in \R$, 
		$$e^x = \limn \bigg(1 + \frac{x}{n} \bigg)^n.$$
	\end{proof}
	
	
	\part{To Do and Ideas}
	
	\chapter{Angle}
	
	Need to talk about the idea of an angle. There are a couple notions here. 
	\begin{itemize}
		\item \tbf{Angle as the surface area measure on the unit circle:} \\
		Here the idea is that angle does not depend on the size of a cirle, is proportional to the surface area of the arc to which it corresponds. Thus we may, by convention define an angle between two (oriented) vectors to be the region of the unit sphere $S^2$ which is bounded by $1$-dimensional geodesic ball on $S^2$ created by the oriented vectors. This idea requires a Banach space and possibly nonsmooth geometry for different geodesic distances like $l_1$ distance. With this, we can define the trig functions by relating the surface area measure to the coordinates and we can prove any number of results about them.
		\item  \tbf{Angle as the measure of orthogonality:} \\
		Here the idea is that in an inner product space, an inner product between unit vectors behaves like the $\cos$ function defined above. In paticular, in the special case of $\R^n$ equipped with the standard inner product, we exactly reproduce the usual notion of angle when the angle is at most $\pi/2$. In addition, the inner product connects the ideas of angle-between and length-of vectors. 
	\end{itemize}
	
	
	
	
	
\end{document}
	
	
	
